\chapter{Turbo, LDPC and Polar Codes}
\label{ch:moderncodes}

\begin{chapterintro}
For forty-five years after Shannon proved that reliable communication was possible at any rate
below capacity, nobody knew how to get there with a decoder that could be built. The codes of
Chapter~\ref{ch:classiccodes} were elegant and useful, but at the error rates that matter they
stayed 2 to 3\,dB from the limit, and every attempt to close the gap with longer codes ran into
decoders whose complexity grew exponentially. Then, within about fifteen years, three families of
codes arrived that closed it almost completely: turbo codes in 1993, low-density parity-check
(LDPC) codes rediscovered in 1996, and polar codes in 2009. They share one idea: a long,
powerful code built from simple pieces, decoded by passing \emph{soft} probabilistic messages
between those pieces. This chapter develops that idea from the log-likelihood ratio upwards. It
derives the BCJR algorithm and builds the turbo decoder around it; explains why turbo codes have
a waterfall and an error floor, and how EXIT charts predict both; develops belief propagation on
the Tanner graph of an LDPC code, its hardware-friendly min-sum approximations, and density
evolution, which computes a code's threshold before a single frame is simulated; shows how
protographs and quasi-cyclic lifting turned LDPC codes into the data channels of 5G, Wi-Fi,
DVB-S2 and 10GBASE-T; and explains channel polarization, successive-cancellation list decoding
and the reasons polar codes carry the control channels of 5G. It closes with the short-blocklength
comparison that decided the 5G coding battle, hybrid ARQ, spatial coupling, the codes of optical
fibre and flash memory, and a look at what might come next.
\end{chapterintro}

%======================================================================
\section{The road to the Shannon limit}
\label{sec:ch15:road}
%======================================================================

Chapter~\ref{ch:infotheory} ended with a promise and a warning. The promise: on the binary-input
additive white Gaussian noise (BI-AWGN) channel, a rate-1/2 code can in principle deliver any
error rate you like at $\ebno=0.19$\,dB, and a rate-1/3 code at $-0.50$\,dB. The warning: the
proof used random codes of enormous length, decoded by comparing the received word with every
codeword. A random code of rate 1/2 and length 1000 has $2^{500}$ codewords; no conceivable
computer will ever decode it that way.

The history of channel coding is the history of the gap between those two statements.
Table~\ref{tab:ch15:gap} gives a rough map. Uncoded BPSK needs about 9.6\,dB to reach a bit error
rate of $10^{-5}$. The first practical codes of the 1950s and 1960s, Hamming and Golay codes
decoded with hard decisions, recovered between a fraction of a decibel and about two decibels. The constraint-length-7 convolutional
code with soft-decision Viterbi decoding (Chapter~\ref{ch:classiccodes}), standardised by NASA in
the 1970s and still inside every Wi-Fi chip, reaches $10^{-5}$ at about 4.4\,dB, 4.2\,dB from the
rate-1/2 limit. Concatenating it with an outer Reed--Solomon code, as the Voyager missions did,
brought the requirement to about 2.5\,dB. The most elaborate schemes of the early 1990s---such
as the constraint-length-15, rate-1/4 code built at JPL for the Galileo mission, whose ``Big Viterbi
Decoder'' tracked 16\,384 states, concatenated with a Reed--Solomon code---were still of the order
of 1.5 to 2\,dB away. Most engineers believed that much of the remaining gap was the price of
practicality. The \term{cutoff rate} $R_0$, the rate up to which sequential decoding works with
bounded average effort, was widely regarded as the ``practical capacity''; on the BPSK/AWGN
channel $R_0=1/2$ requires $\ebno\approx2.5$\,dB, about 2.3\,dB more than capacity, and at very
low rates the gap grows to 3\,dB.

\begin{table}[htbp]\centering\small
\caption{Approximate $\ebno$ required for a bit error rate of $10^{-5}$ on the AWGN channel with
BPSK or QPSK, for landmark coding schemes. Values are typical published figures, rounded; the
BPSK-input Shannon limits are $0.19$\,dB at rate 1/2 and $-0.50$\,dB at rate 1/3.}
\label{tab:ch15:gap}
\begin{tabularx}{\linewidth}{@{}l l c c X@{}}
\toprule
Scheme & Year & Rate & $\ebno$ (dB) & Remark\\
\midrule
Uncoded BPSK & -- & 1 & 9.6 & the reference\\
Hamming (7,4), hard decisions & 1950 & 0.57 & $\approx$9.2 & less than half a decibel of gain\\
Convolutional $K=7$, soft Viterbi & 1970s & 1/2 & $\approx$4.4 & NASA standard; 802.11a/g\\
RS(255,223) + $K=7$ convolutional & 1977 & $\approx$0.44 & $\approx$2.5 & Voyager, CCSDS\\
Turbo code, $K=65\,536$, 18 iterations & 1993 & 1/2 & 0.7 & Berrou \emph{et al.}; 0.5\,dB from the limit\\
Irregular LDPC, $n=10^7$ & 2001 & 1/2 & $\approx$0.23 & Chung \emph{et al.}, $10^{-6}$; 0.04\,dB from the limit\\
5G NR LDPC, a few thousand bits & 2017 & 1/2 & $\approx$1--1.5 & practical blocks at BLER $10^{-2}$\\
\bottomrule
\end{tabularx}
\end{table}

The breakthrough, when it came, did not come from a better algebraic construction. It came from
two changes of viewpoint. First, a decoder should not decide bits; it should compute
\emph{probabilities} of bits and pass them on. Second, a long code does not need a
long-code decoder: if the code is built from small component codes, each with a cheap
probabilistic decoder, the components can decode \emph{each other}, iteratively, by exchanging
probabilities. Both ideas were available in 1963, in Robert Gallager's thesis on low-density
parity-check codes. Both were ignored for thirty years.

\begin{history}[Geneva, May 1993]
At the IEEE International Conference on Communications in Geneva in May 1993, Claude Berrou, Alain
Glavieux and Punya Thitimajshima of the \'Ecole Nationale Sup\'erieure des T\'el\'ecommunications
de Bretagne presented a paper titled ``Near Shannon limit error-correcting coding and decoding:
Turbo-codes''. Their rate-1/2 code, two 16-state recursive convolutional encoders joined by a
$256\times256$ interleaver, reached a bit error rate of $10^{-5}$ at $\ebno=0.7$\,dB after 18
iterations of decoding, about half a decibel from the Shannon limit. The claim was more than a
decibel better than anything previously built, and many in the audience assumed a mistake:
perhaps an error in the SNR normalisation, perhaps a factor of two between $E_b$ and $E_s$.
Berrou and Glavieux were not coding theorists by background---Berrou was an electronics engineer
who had been designing Viterbi decoder chips---and the paper's explanation of \emph{why} it worked
was sketchy. Within months, however, groups at JPL and elsewhere had reproduced the results. The
name came from the turbocharger: the decoder's output is fed back to its input, as a turbocharger
feeds exhaust back to drive the intake. Seven years later, Berrou, Glavieux and Thitimajshima's
paper had become one of the most cited in communications, and turbo codes were written into the
third-generation mobile standard and NASA's deep-space recommendations.
\end{history}

The rest of this chapter follows the ideas in logical rather than strict historical order. We
start with the currency that every modern decoder trades in, the log-likelihood ratio, and the
algorithm that computes exact bit probabilities on a trellis.

%======================================================================
\section{Soft information and the BCJR algorithm}
\label{sec:ch15:soft}
%======================================================================

\subsection{Log-likelihood ratios}

Chapter~\ref{ch:modulation} introduced the \term{log-likelihood ratio} (LLR) of a bit as the soft
output of a demodulator. We now make it the central object. For a binary random variable
$c\in\{0,1\}$ and an observation $y$, define
\begin{equation}
L(c\,|\,y)=\ln\frac{P(c=0\,|\,y)}{P(c=1\,|\,y)}.
\label{eq:ch15:llr}
\end{equation}
The sign is the hard decision (positive means 0, with the BPSK mapping $0\mapsto+1$,
$1\mapsto-1$ used throughout the book and in \texttt{commlib}), and the magnitude is the
confidence: $|L|=0$ means ``no idea'', $|L|=4.6$ means the decision is wrong with probability
$1/(1+e^{4.6})\approx1\%$, and $|L|=\infty$ means certainty. Probabilities are recovered as
\begin{equation}
P(c=0\,|\,y)=\frac{1}{1+e^{-L}},\qquad P(c=1\,|\,y)=\frac{1}{1+e^{L}}.
\end{equation}
By Bayes' rule the posterior LLR splits into two terms,
\begin{equation}
L(c\,|\,y)=\underbrace{\ln\frac{p(y\,|\,c=0)}{p(y\,|\,c=1)}}_{L_{\text{ch}}}+
\underbrace{\ln\frac{P(c=0)}{P(c=1)}}_{L_A},
\label{eq:ch15:bayes}
\end{equation}
the \emph{channel} LLR and the \emph{a-priori} LLR. For BPSK in real Gaussian noise of variance
$\sigma^2$, $y=x+n$ with $x=1-2c$, the channel LLR is linear in the observation:
\begin{equation}
L_{\text{ch}}=\ln\frac{e^{-(y-1)^2/2\sigma^2}}{e^{-(y+1)^2/2\sigma^2}}=\frac{2}{\sigma^2}\,y.
\label{eq:ch15:chllr}
\end{equation}
A receiver that scales its matched-filter samples by $2/\sigma^2$ therefore already produces
LLRs. Given $x=+1$, $L_{\text{ch}}$ is Gaussian with mean $\mu=2/\sigma^2$ and variance
$2\mu$. That relationship---variance equal to twice the mean---is the \term{consistency
condition}, and it holds for any correctly computed LLR whose distribution is Gaussian; we will
use it in density evolution and EXIT charts.

Two properties make LLRs the natural currency of iterative decoding. First,
\emph{independent} evidence about the same bit adds: if $y_1$ and $y_2$ are conditionally
independent given $c$, $L(c\,|\,y_1,y_2)=L(c\,|\,y_1)+L(c\,|\,y_2)$. That is the whole of a
repetition code's decoder and of the variable-node update of belief propagation. Second, the LLR
of the modulo-2 sum of two independent bits has a simple closed form. If $c=c_1\oplus c_2$,
\begin{equation}
L(c_1\oplus c_2)=2\tanh^{-1}\!\Big(\tanh\frac{L_1}{2}\tanh\frac{L_2}{2}\Big)
\;\triangleq\;L_1\boxplus L_2,
\label{eq:ch15:boxplus}
\end{equation}
which follows from $P(c=0)-P(c=1)=\tanh(L/2)$ and the fact that the ``bias'' $P(0)-P(1)$ of an XOR
of independent bits is the product of their biases. The operation $\boxplus$, the
\term{box-plus} or check-node operation, is the whole of a single parity check's decoder. Its
magnitude is always smaller than the smaller input magnitude, and a good approximation is
\begin{equation}
L_1\boxplus L_2\approx\operatorname{sign}(L_1)\operatorname{sign}(L_2)\min(|L_1|,|L_2|),
\label{eq:ch15:minapprox}
\end{equation}
since an XOR is only as reliable as its least reliable input. Everything in this chapter's
decoders is built from sums, box-pluses, and their approximations.

\begin{worked}[Box-plus arithmetic]
Two bits have LLRs $L_1=+2.0$ and $L_2=-1.0$. What is known about $c_1\oplus c_2$?

$\tanh(1.0)=0.762$ and $\tanh(-0.5)=-0.462$, so the product is $-0.352$ and
$L_1\boxplus L_2=2\tanh^{-1}(-0.352)=-0.735$. The XOR is probably 1 (as the hard decisions
$0\oplus1$ suggest) but with less confidence than either input. The min-sum
approximation~\eqref{eq:ch15:minapprox} gives $-1.0$, overestimating the confidence by about a
third. The exact correction is
$L_1\boxplus L_2=\operatorname{sign}\cdot\min+\ln(1+e^{-|L_1+L_2|})-\ln(1+e^{-|L_1-L_2|})
=-1.0+0.313-0.049=-0.735$: a small table lookup restores the exact value.
\end{worked}

\subsection{A-priori, extrinsic and a-posteriori information}

Consider a decoder for one component code. For each information bit $u_k$ it receives a channel
LLR $L_{\text{ch},k}$ (from the systematic bit), possibly an a-priori LLR $L_{A,k}$ (from somewhere
else---another decoder), and channel LLRs for the parity bits. It computes the a-posteriori LLR
$L(u_k\,|\,\text{everything})$. It can be shown (and we will see it concretely in the BCJR
algorithm) that for a systematic code this output always decomposes as
\begin{equation}
L_{\text{APP},k}=L_{\text{ch},k}+L_{A,k}+L_{E,k}.
\label{eq:ch15:extrinsic}
\end{equation}
The third term, the \term{extrinsic information}, is what the code's structure---the \emph{other}
bits, through the parity constraints---says about $u_k$. It is the only part of the output that
is new to whoever supplied $L_{A,k}$. The central rule of iterative decoding follows: \emph{pass
on only extrinsic information}. If a decoder passed the full APP back to the decoder that gave it
the a-priori value, that decoder would be told its own opinion as if it were independent
evidence, count it twice, and become confidently wrong. Subtracting what the receiver already knew
is the discipline that makes iterative decoding approximately Bayesian.

\begin{keyidea}[Soft in, soft out]
A \term{soft-input soft-output} (SISO) decoder takes LLRs in and gives LLRs out. Two SISO decoders
of two different codes on the same bits can help each other: each turns the other's extrinsic
output into its own a-priori input. Every capacity-approaching decoder of this chapter---turbo,
LDPC and, in a sequential form, polar---is a network of simple SISO decoders exchanging extrinsic
LLRs.
\end{keyidea}

\subsection{The trellis and the MAP decoder}

The Viterbi algorithm of Chapter~\ref{ch:classiccodes} finds the most likely \emph{sequence} of
states through a trellis. It minimises the probability of a sequence error but outputs only hard
decisions. The \term{BCJR algorithm}, published by Lalit Bahl, John Cocke, Frederick Jelinek and
Josef Raviv in 1974, computes instead the a-posteriori probability of each \emph{individual}
bit, which is exactly the soft output a SISO decoder needs.

Let the encoder be a finite-state machine with states $s\in\{0,\dots,S-1\}$. At time $k$ an input
bit $u_k$ drives it from state $s'$ to state $s$ and produces output bits whose noisy observations
we collect in $\vect y_k$. The whole received sequence is $\vect y=(\vect y_1,\dots,\vect y_N)$. The
APP LLR of $u_k$ is a ratio of sums over trellis branches:
\begin{equation}
L(u_k\,|\,\vect y)=\ln\frac{\sum_{(s',s):\,u_k=0}p(s_{k-1}=s',s_k=s,\vect y)}
{\sum_{(s',s):\,u_k=1}p(s_{k-1}=s',s_k=s,\vect y)}.
\label{eq:ch15:app}
\end{equation}
Because the channel is memoryless and the encoder is Markov, the joint probability of a branch and
the observations factors into three pieces: everything before the branch, the branch itself, and
everything after it,
\begin{equation}
p(s',s,\vect y)=\alpha_{k-1}(s')\,\gamma_k(s',s)\,\beta_k(s),
\label{eq:ch15:factor}
\end{equation}
with
\begin{align}
\alpha_k(s)&=p(s_k=s,\vect y_1,\dots,\vect y_k), \nonumber\\
\gamma_k(s',s)&=p(s_k=s,\vect y_k\,|\,s_{k-1}=s') = P(u_k)\,p(\vect y_k\,|\,s',s), \\
\beta_k(s)&=p(\vect y_{k+1},\dots,\vect y_N\,|\,s_k=s). \nonumber
\end{align}
The forward and backward quantities obey simple recursions,
\begin{equation}
\alpha_k(s)=\sum_{s'}\alpha_{k-1}(s')\,\gamma_k(s',s),\qquad
\beta_{k-1}(s')=\sum_{s}\gamma_k(s',s)\,\beta_k(s),
\label{eq:ch15:recursions}
\end{equation}
started from the known initial state ($\alpha_0(0)=1$) and, for a terminated trellis, the known final
state ($\beta_N(0)=1$). The algorithm is therefore: compute all branch metrics $\gamma$; run a forward
pass for $\alpha$; run a backward pass for $\beta$; and combine them with~\eqref{eq:ch15:app}. It
costs about three times a Viterbi decoder, and it is exact.

\begin{figure}[htbp]\centering
\begin{tikzpicture}[x=1.55cm,y=0.62cm,
  st/.style={circle,draw=navy,fill=navy!8,thick,inner sep=0pt,minimum size=4.2mm,font=\sffamily\tiny},
  b0/.style={navy!55,thick}, b1/.style={accent!70,thick,dashed}]
\foreach \k in {0,...,5}{
  \foreach \s in {0,...,3}{ \node[st] (n\k\s) at (\k,-\s) {\s}; } }
\foreach \k [evaluate=\k as \kk using int(\k+1)] in {0,...,4}{
  \draw[b0] (n\k0) -- (n\kk0); \draw[b1] (n\k0) -- (n\kk2);
  \draw[b0] (n\k1) -- (n\kk2); \draw[b1] (n\k1) -- (n\kk0);
  \draw[b0] (n\k2) -- (n\kk1); \draw[b1] (n\k2) -- (n\kk3);
  \draw[b0] (n\k3) -- (n\kk3); \draw[b1] (n\k3) -- (n\kk1); }
\foreach \k in {0,...,5}{ \node[font=\sffamily\scriptsize,navy] at (\k,0.85) {$k=\k$}; }
\fill[navy!10,rounded corners] (-0.35,-4.05) rectangle (2.3,-4.85);
\node[font=\sffamily\scriptsize,navy] at (0.98,-4.45) {$\alpha_{k-1}(s')$: past};
\fill[accent!10,rounded corners] (2.35,-4.05) rectangle (3.65,-4.85);
\node[font=\sffamily\scriptsize,accent] at (3.0,-4.45) {$\gamma_k(s',s)$};
\fill[practc!12,rounded corners] (3.7,-4.05) rectangle (5.35,-4.85);
\node[font=\sffamily\scriptsize,practc] at (4.52,-4.45) {$\beta_k(s)$: future};
\draw[-{Stealth[length=2mm]},navy,very thick] (0,1.5) -- (2,1.5) node[midway,above,font=\sffamily\scriptsize]{forward recursion};
\draw[-{Stealth[length=2mm]},practc,very thick] (5,1.5) -- (3,1.5) node[midway,above,font=\sffamily\scriptsize]{backward recursion};
\draw[line width=2.2pt,accent!80,opacity=0.5] (n21) -- (n30);
\node[font=\sffamily\scriptsize,anchor=west] at (5.3,-0.5) {\textcolor{navy!70}{\rule[2pt]{5mm}{1pt}} $u_k=0$};
\node[font=\sffamily\scriptsize,anchor=west] at (5.3,-1.3) {\textcolor{accent!80}{- - -} $u_k=1$};
\end{tikzpicture}
\caption{The BCJR algorithm on a 4-state trellis. The probability of any branch---here the
highlighted one from state 1 at time 2 to state 0 at time 3---given all the observations is the
product of a forward metric $\alpha$ (summarising the observations before the branch), the branch
metric $\gamma$ and a backward metric $\beta$ (summarising the observations after it). Summing
over all branches labelled $u_k=0$ and all branches labelled $u_k=1$ gives the bit's APP LLR.}
\label{fig:ch15_bcjr_trellis}
\end{figure}

\subsection{The log domain: log-MAP and max-log-MAP}

Products of probabilities underflow after a few hundred steps, so practical decoders work with
logarithms: $A_k(s)=\ln\alpha_k(s)$, $B_k(s)=\ln\beta_k(s)$, $\Gamma_k=\ln\gamma_k$. Products
become sums, and a sum of probabilities becomes the \term{Jacobian logarithm}
\begin{equation}
\max{}^*(a,b)\;\triangleq\;\ln(e^a+e^b)=\max(a,b)+\ln\!\big(1+e^{-|a-b|}\big).
\label{eq:ch15:maxstar}
\end{equation}
The recursions become
\begin{equation}
A_k(s)=\max_{s'}{}^*\big[A_{k-1}(s')+\Gamma_k(s',s)\big],\qquad
B_{k-1}(s')=\max_{s}{}^*\big[\Gamma_k(s',s)+B_k(s)\big],
\end{equation}
and the output is
\begin{equation}
L(u_k)=\max_{u_k=0}{}^*\big[A_{k-1}(s')+\Gamma_k(s',s)+B_k(s)\big]-\max_{u_k=1}{}^*\big[A_{k-1}(s')+\Gamma_k(s',s)+B_k(s)\big].
\end{equation}
This is the \term{log-MAP} algorithm; it is exactly equivalent to BCJR. The correction term
$\ln(1+e^{-|a-b|})$ is at most $\ln2\approx0.69$ and falls below 0.05 when $|a-b|>3$, so hardware
implements it with a lookup table of eight or so entries. Dropping it altogether gives
\term{max-log-MAP}, in which every $\max^*$ becomes a plain $\max$. The forward and backward passes
are then two Viterbi algorithms running in opposite directions, and the output for each bit is the
difference between the best path with $u_k=0$ and the best path with $u_k=1$. Max-log-MAP loses
about 0.3 to 0.5\,dB in a turbo decoder, but it is insensitive to the scaling of the input LLRs
(an unknown noise variance does no harm) and most of its loss is recovered by multiplying the
extrinsic output by a factor of about 0.7, a trick found empirically by several groups around
2000 (Figure~\ref{fig:ch15_turbo_iters}).

With BPSK on the AWGN channel and the LLR form of the inputs, the branch metric of a rate-1/2
systematic code whose branch $(s',s)$ carries input $u$ and parity $p$, in bipolar form
$x^u=1-2u$ and $x^p=1-2p$, is
\begin{equation}
\Gamma_k(s',s)=\tfrac12x^u\big(L_{\text{ch},k}^{u}+L_{A,k}\big)+\tfrac12x^p\,L_{\text{ch},k}^{p}
\label{eq:ch15:branch}
\end{equation}
(up to a constant that cancels). Substituting into the output shows that the terms involving
$L^u_{\text{ch},k}+L_{A,k}$ factor out of both sums, leaving exactly the
decomposition~\eqref{eq:ch15:extrinsic}: the extrinsic output is computed from the parity
observations and from everything the trellis says about the \emph{other} bits.

\begin{history}[A twenty-year sleeper]
The BCJR paper appeared in the \emph{IEEE Transactions on Information Theory} in March 1974, from
IBM's Thomas J.~Watson Research Center, where Bahl, Jelinek and Raviv were working on speech
recognition. (The same forward--backward recursions are the Baum--Welch algorithm of hidden Markov
models.) For decoding, it was judged a curiosity: it gave almost exactly the same bit error rate as
the Viterbi algorithm at roughly three times the cost. What nobody needed in 1974 was a decoder
whose \emph{output} was soft. When Berrou's turbo decoder needed exactly that in 1993, the BCJR
algorithm was waiting, and it has been the engine of every turbo decoder since. Hagenauer and
Hoeher's soft-output Viterbi algorithm (SOVA, 1989) is a cheaper approximation that was also used in
early turbo decoders.
\end{history}

%======================================================================
\section{Turbo codes}
\label{sec:ch15:turbo}
%======================================================================

\subsection{Parallel concatenation}

A \term{turbo code} (more precisely a \term{parallel concatenated convolutional code}, PCCC)
encodes a block of $K$ information bits three times: once as they are (the systematic bits), once
with a convolutional encoder, and once with a second convolutional encoder that sees the same bits
in a permuted order (Figure~\ref{fig:ch15_turbo_encoder}). The permutation, the
\term{interleaver} $\pi$, is the heart of the code. The rate is 1/3 before puncturing; deleting
some parity bits gives any higher rate.

\begin{figure}[htbp]\centering
\begin{tikzpicture}[node distance=0.5cm and 0.9cm,
  blk/.style={draw=navy,fill=navy!6,thick,minimum height=0.9cm,minimum width=1.9cm,align=center,font=\sffamily\scriptsize},
  arr/.style={-{Stealth[length=2mm]},thick,navy},
  reg/.style={draw=navy,thick,minimum size=0.55cm,font=\sffamily\scriptsize},
  xo/.style={circle,draw=navy,thick,inner sep=0pt,minimum size=3.4mm}]
\node[font=\sffamily\scriptsize,align=center] (in) at (0,0) {$u_k$\\$K$ bits};
\node[blk,right=1.6cm of in] (e1) {RSC encoder 1\\$(13,15)_8$};
\node[blk,below=0.9cm of e1,fill=accent!6,draw=accent] (pi) {interleaver $\pi$};
\node[blk,right=0.9cm of pi] (e2) {RSC encoder 2\\$(13,15)_8$};
\node[font=\sffamily\scriptsize,anchor=west] (xs) at (7.2,0.9) {$x_k=u_k$ (systematic)};
\node[font=\sffamily\scriptsize,anchor=west] (z1) at (7.2,0) {$z_k$ (parity 1)};
\node[font=\sffamily\scriptsize,anchor=west] (z2) at (7.2,-1.4) {$z'_k$ (parity 2)};
\draw[arr] (in) -- (e1);
\coordinate (sp) at (1.0,0);
\fill[navy] (sp) circle (1.5pt);
\draw[arr] (sp) |- (xs.west);
\draw[arr] (sp) |- (pi.west);
\draw[arr] (e1) -- (z1.west);
\draw[arr] (pi) -- (e2);
\draw[arr] (e2) -- (z2.west -| e2.east) -- (z2.west);
% RSC detail
\begin{scope}[shift={(0.6,-3.6)}]
\node[font=\sffamily\scriptsize,anchor=west,navy] at (-0.55,0.95) {Inside each RSC encoder (8 states, TS 36.212):};
\node[xo] (a1) at (0,0) {\tiny$+$};
\node[reg] (d1) at (1.3,0) {$D$};
\node[reg] (d2) at (2.3,0) {$D$};
\node[reg] (d3) at (3.3,0) {$D$};
\node[xo] (fb) at (0,-0.9) {\tiny$+$};
\node[xo] (o1) at (4.6,0.55) {\tiny$+$};
\node[font=\sffamily\scriptsize] (uin) at (-1.1,0) {$u_k$};
\draw[arr] (uin) -- (a1);
\draw[arr] (a1) -- node[above,font=\tiny]{$a_k$} (d1); \draw[arr] (d1) -- (d2); \draw[arr] (d2) -- (d3);
\draw[arr] (d2.south) |- (fb.east);
\draw[arr] (d3.south) -- ++(0,-0.9) -- (fb.east);
\draw[arr] (fb) -- (a1);
\draw[arr] (0.6,0) |- (o1.west);
\draw[arr] (d1.north) |- ([yshift=1mm]o1.west);
\draw[arr] (d3.north) |- ([yshift=-1mm]o1.west);
\node[font=\sffamily\scriptsize,anchor=west] (zo) at (5.1,0.55) {$z_k$};
\draw[arr] (o1) -- (zo);
\node[font=\sffamily\tiny,anchor=west,align=left] at (4.1,-0.7) {feedback $1+D^2+D^3$ (13 octal)\\feed-forward $1+D+D^3$ (15 octal)};
\end{scope}
\end{tikzpicture}
\caption{The turbo encoder of UMTS and LTE (3GPP TS 25.212 and TS 36.212). Two identical
eight-state recursive systematic convolutional (RSC) encoders see the information bits in natural
and in interleaved order. Only one copy of the systematic bits is sent, so the rate is 1/3, plus
twelve tail bits that return both encoders to the zero state. Below: one constituent encoder; the
feedback makes its impulse response infinite.}
\label{fig:ch15_turbo_encoder}
\end{figure}

\subsection{Why the constituent codes must be recursive}

The constituent encoders are \term{recursive systematic convolutional} (RSC) codes: the parity is
computed from a register driven by the input \emph{plus feedback} from its own contents. The LTE
code has feedback polynomial $g_0(D)=1+D^2+D^3$ and feed-forward polynomial $g_1(D)=1+D+D^3$; its
transfer function is $z(D)=u(D)\,g_1(D)/g_0(D)$. The feedback is not a detail. It is what gives
turbo codes their power, and it is worth understanding exactly why.

A non-recursive (feed-forward) encoder answers a weight-1 input with a finite, low-weight output:
a single 1 produces one copy of the generator. If the constituent codes were feed-forward, a single
information 1 would produce low-weight parity in \emph{both} encoders, wherever the interleaver
moved it, and the code would have codewords of weight about $1+2\cdot 3$ no matter how long the
block. A recursive encoder answers a weight-1 input with an \emph{infinite} impulse response: the
parity keeps toggling until the end of the block. Only inputs whose polynomial $u(D)$ is divisible
by $g_0(D)$ return the encoder to the zero state and produce low-weight parity. Because
$1+D^2+D^3$ is primitive of degree 3, the lowest-weight such inputs are pairs of ones separated by
a multiple of $2^3-1=7$ positions: $u(D)=D^i(1+D^7)$ gives a parity sequence of weight 6,
$D^i(1+D^{14})$ weight 10, and in general $1+D^{7t}$ weight $4t+2$.

Now the interleaver enters. A weight-2 input pattern that is ``bad'' for encoder 1 (separation 7)
is scattered by $\pi$ to two positions whose separation is, for a good interleaver, almost never a
multiple of 7 and almost never small; encoder 2 then produces high-weight parity. A low-weight
codeword needs an input pattern that is bad for \emph{both} encoders simultaneously, and for a
random interleaver of size $K$ the probability that a given weight-2 pattern survives is of order
$1/K$. This is the \term{interleaver gain}: Benedetto and Montorsi showed in 1996 that for
recursive constituent codes the bit error probability in the high-SNR region falls as $K^{-1}$ as
the interleaver grows, whereas for feed-forward constituents there is no such gain. A turbo code
does not have a large minimum distance---its $d_{\min}$ grows only slowly (at most
logarithmically) with $K$---but it has very \emph{few} low-weight codewords, a property sometimes
called \term{spectral thinning}. At low SNR, where a code's performance is set by the enormous
number of codewords at moderate distance rather than by the few at the minimum, that is what
matters.

\subsection{The interleaver}

For short blocks, the interleaver design matters as much as the constituent codes. Its jobs are to
break the low-weight input patterns described above and to make the two decoders' views of the
data as independent as possible, so that the extrinsic information one passes to the other is
genuinely new. Designs have progressed through several generations:
\begin{itemize}
\item \textbf{Random} interleavers work well on average for long blocks but occasionally map a
short pattern to a short pattern.
\item \textbf{S-random} interleavers (Divsalar and Pollara, 1995) are random permutations with the
constraint that any two inputs within $S$ positions of each other are mapped at least $S$ apart. A
spread of $S\approx\sqrt{K/2}$ is achievable and removes the worst weight-2 patterns.
\item \textbf{Algebraic} interleavers can be computed on the fly instead of stored, and can be
designed to be \emph{contention-free}: a decoder that splits the block into $P$ windows processed by
$P$ parallel BCJR engines can read and write its extrinsic memory in $P$ banks without two engines
ever needing the same bank in the same clock cycle. That property is essential at LTE data rates.
\end{itemize}
The UMTS turbo code used a prime-number-based interleaver defined by a multi-step procedure over a
row--column array. LTE replaced it with the \term{quadratic permutation polynomial} (QPP)
interleaver of Sun and Takeshita (2005),
\begin{equation}
\pi(i)=\big(f_1\,i+f_2\,i^2\big)\bmod K,\qquad i=0,\dots,K-1,
\label{eq:ch15:qpp}
\end{equation}
with $(f_1,f_2)$ tabulated in TS 36.212 for each of the 188 allowed block sizes from $K=40$ to
$K=6144$. A QPP is a permutation when $f_1$ is coprime to $K$ and every prime factor of $K$ divides
$f_2$. It is contention-free for every window size that divides $K$, it needs no memory, and it can
be generated recursively with additions only, since
$\pi(i+1)=\pi(i)+g(i)$ with $g(i+1)=g(i)+2f_2 \bmod K$.

\begin{worked}[The QPP interleaver for $K=40$]
The smallest LTE block uses $K=40$, $f_1=3$, $f_2=10$. Check that it is a permutation and
compute the first few entries.

$K=40=2^3\cdot5$; $f_1=3$ is coprime to 40, and $f_2=10=2\cdot5$ contains both prime factors of
$K$, so~\eqref{eq:ch15:qpp} is a permutation. The first entries are $\pi(0)=0$,
$\pi(1)=13$, $\pi(2)=(6+40)\bmod40=6$, $\pi(3)=(9+90)\bmod40=19$, $\pi(4)=(12+160)\bmod40=12$,
$\pi(5)=(15+250)\bmod 40=25$. With the recursion: $g(0)=f_1+f_2=13$, and $g$ increases by
$2f_2=20$ each step (mod 40), so $\pi$ advances by $13,33,13,33,\dots$ (mod 40). Two inputs
$d$ positions apart are mapped to positions that differ by $f_1d+f_2d(2i+d)$, which varies with
$i$---that is what breaks up regular patterns. The last LTE size, $K=6144$, uses $f_1=263$,
$f_2=480$.
\end{worked}

\subsection{Iterative decoding}

The optimal decoder for a turbo code would compute APPs on the joint trellis of both encoders and
the interleaver, a machine with astronomically many states. The \term{turbo decoder}
(Figure~\ref{fig:ch15_turbo_decoder}) instead uses two BCJR decoders, one per constituent code, and
lets them exchange extrinsic information:
\begin{enumerate}
\item Decoder 1 receives the systematic LLRs $L_{\text{ch}}^{x}$, the parity-1 LLRs
$L_{\text{ch}}^{z}$ and an a-priori input $L_{A1}$ (zero in the first iteration). It runs BCJR and
outputs extrinsic LLRs $L_{E1}=L_{\text{APP}1}-L_{\text{ch}}^{x}-L_{A1}$.
\item The extrinsic values are interleaved and become decoder 2's a-priori input,
$L_{A2}=\pi(L_{E1})$. Decoder 2 uses the interleaved systematic LLRs and the parity-2 LLRs and
outputs $L_{E2}=L_{\text{APP}2}-\pi(L_{\text{ch}}^{x})-L_{A2}$.
\item These are de-interleaved and become decoder 1's new a-priori input,
$L_{A1}=\pi^{-1}(L_{E2})$. One pass through both decoders is one \emph{iteration}.
\item After a fixed number of iterations, or when a CRC passes or the decisions stop changing,
decide $\hat u_k=0$ if $L_{\text{ch},k}^{x}+L_{A1,k}+L_{E1,k}>0$.
\end{enumerate}

\begin{figure}[htbp]\centering
\begin{tikzpicture}[node distance=0.6cm and 0.8cm,
  blk/.style={draw=navy,fill=navy!6,thick,minimum height=1.25cm,minimum width=2.1cm,align=center,font=\sffamily\scriptsize},
  il/.style={draw=accent,fill=accent!6,thick,minimum height=0.7cm,minimum width=0.9cm,font=\sffamily\scriptsize},
  arr/.style={-{Stealth[length=2mm]},thick,navy},
  sm/.style={circle,draw=navy,thick,inner sep=0pt,minimum size=4mm,font=\scriptsize}]
\node[blk] (d1) at (0,0) {SISO decoder 1\\(BCJR, code 1)};
\node[blk] (d2) at (6.2,0) {SISO decoder 2\\(BCJR, code 2)};
\node[sm] (s1) at (2.1,0.2) {$-$};
\node[il] (pi1) at (3.4,0.2) {$\pi$};
\node[sm] (s2) at (8.3,0.2) {$-$};
\node[il] (pi2) at (3.4,-1.7) {$\pi^{-1}$};
\node[il] (pi3) at (3.4,1.55) {$\pi$};
\node[font=\sffamily\scriptsize,anchor=east] (ys) at (-2.1,0.35) {$L^{x}_{\text{ch}}$};
\node[font=\sffamily\scriptsize,anchor=east] (yp1) at (-2.1,-0.35) {$L^{z}_{\text{ch}}$};
\node[font=\sffamily\scriptsize,anchor=west] (yp2) at (6.2,-1.05) {$L^{z'}_{\text{ch}}$};
\draw[arr] (ys) -- ([yshift=3.5mm]d1.west);
\draw[arr] (yp1) -- ([yshift=-3.5mm]d1.west);
\draw[arr] ([yshift=2mm]d1.east) -- node[above,font=\tiny]{$L_{\text{APP}1}$} (s1);
\draw[arr] (s1) -- node[above,font=\tiny]{$L_{E1}$} (pi1);
\draw[arr] (pi1) -- node[above,font=\tiny]{$L_{A2}$} ([yshift=2mm]d2.west);
\draw[arr] (-1.6,0.35) |- (pi3.west);
\draw[arr] (pi3.east) -| ([xshift=-4mm]d2.north west) -- ++(0,-0.3);
\draw[arr] (yp2) -- (d2.south);
\draw[arr] ([yshift=2mm]d2.east) -- node[above,font=\tiny]{$L_{\text{APP}2}$} (s2);
\draw[arr] (s2) -- ++(0,-1.9) -| ++(0,0) -- (pi2.east) node[pos=0.3,above,font=\tiny]{$L_{E2}$};
\draw[arr] (pi2.west) -| ([xshift=4mm]d1.south west) node[pos=0.2,above,font=\tiny]{$L_{A1}$};
\draw[arr,accent] (d2.north) ++(0.6,0) -- ++(0,0.9) node[above,font=\sffamily\scriptsize]{decisions $\hat u$ (after $\pi^{-1}$)};
\draw[navy!60] (2.1,1.0) -- (s1) node[pos=0,above,font=\tiny]{$L^x_{\text{ch}}+L_{A1}$};
\draw[navy!60] (8.3,1.0) -- (s2) node[pos=0,above,font=\tiny]{$\pi(L^x_{\text{ch}})+L_{A2}$};
\end{tikzpicture}
\caption{The iterative turbo decoder. Each SISO decoder subtracts its inputs from its
a-posteriori output so that only extrinsic information is passed on; interleavers keep the two
decoders' views of the bits aligned. The loop is what gave turbo codes their name.}
\label{fig:ch15_turbo_decoder}
\end{figure}

Figure~\ref{fig:ch15_turbo_iters} shows the decoder at work on the $K=1024$ LTE code. After one
iteration the code is only about as good as a single convolutional code. Each further iteration
moves the curve to the left, by a lot at first and then by less, until after about eight iterations
nothing more is gained; the final curve drops from $10^{-1}$ to $10^{-5}$ in about a decibel. The
iterations do not converge to the maximum-likelihood decision---iterative decoding on a graph with
cycles is an approximation---but in the waterfall region they come remarkably close to it.

\mdcfig{ch15_turbo_iters}{Bit error rate of the rate-1/3 LTE turbo code with $K=1024$ (QPP
interleaver $f_1=31$, $f_2=64$, both encoders terminated, $n=3084$) on the BPSK/AWGN channel, after
1 to 8 iterations of log-MAP decoding. Also shown: 8 iterations of max-log-MAP, without and with
the extrinsic scaled by 0.7. Monte Carlo with 200-frame batches, up to 1200 frames per point.}

\begin{pitfall}[Double counting]
The most common bug in a first turbo decoder is to pass the APP output, or to forget to subtract
the a-priori input, so that each decoder is fed its own opinion back. The decoder then converges
fast, confidently and often to the wrong answer; the BER curve looks like an uncoded curve shifted
by a decibel or two, with no waterfall. The second most common bug is a mismatched LLR sign
convention between the demapper and the decoder, which produces a BER near 0.5 or, more
insidiously, one that improves with SNR only because the systematic bits are good. Always check a
new decoder at high SNR with the a-priori input disabled: it must then behave as a single
convolutional decoder.
\end{pitfall}

\subsection{The waterfall and the error floor}

Every turbo code's error-rate curve has two regions (Figures~\ref{fig:ch15_turbo_length}
and~\ref{fig:ch15_turbo_floor}). In the \term{waterfall} region, just above a threshold SNR, the
error rate plunges by several decades per decibel: the iterative decoder either converges (and the
frame is decoded with essentially no errors) or it does not. The position of the waterfall is set
by the \emph{convergence} of the iterative decoder, which depends on the constituent codes and the
code rate but hardly at all on the interleaver; its steepness grows with the block length, because
long blocks average the noise better. At high SNR the curve flattens into an \term{error floor},
where the rare errors that remain are caused by the code's few low-weight codewords. Those are
maximum-likelihood errors: even an ideal decoder would make them. Their contribution is well
approximated by the first terms of the union bound,
\begin{equation}
P_b\approx\sum_{d}\frac{w_d}{K}\,A_d\,\Q\!\Big(\sqrt{2Rd\,\ebno}\Big),
\label{eq:ch15:floor}
\end{equation}
where $A_d$ is the number of codewords of weight $d$ and $w_d$ their average information weight.
For a turbo code the dominant terms come from the weight-2 input patterns of the previous
subsections, and the factor $w_d/K$ is the interleaver gain.

\mdcfig{ch15_turbo_length}{Frame error rate of LTE-type turbo codes (rate 1/3, QPP interleavers
of the TS 36.212 form, 8 log-MAP iterations) for four block lengths. Stars mark the
normal-approximation bound for the best possible code of the same length and rate at FER
$10^{-2}$ (Section~\ref{sec:ch15:compare}). Monte Carlo: at least 60 frame errors, or 300 to 3000
frames, per point.}

Figure~\ref{fig:ch15_turbo_length} shows the effect of block length. At $K=40$ the curve is
shallow and reaches a frame error rate of $10^{-2}$ only at about 2.7\,dB, more than 3\,dB to the
right of the rate-1/3 limit; at $K=6144$ it falls from 1 to $10^{-2}$ within about 0.4\,dB and
reaches it at roughly 0.4\,dB, within a decibel of the limit. The stars show the best that
\emph{any} code of each length and rate could do, according to the normal approximation of
Chapter~\ref{ch:infotheory}. About half of the short code's distance from the limit is that
unavoidable finite-length penalty; the turbo codes themselves sit about 0.6\,dB (long blocks) to
1.2\,dB (short blocks) from the stars.

\mdcfig{ch15_turbo_floor}{The error floor. Bit error rate of two rate-1/3 turbo codes with $K=1024$
that differ only in the interleaver: the QPP interleaver ($f_1=31$, $f_2=64$) and a random
permutation (8 log-MAP iterations, up to 1200 frames per point; points with no errors are not
plotted). Dashed: the contribution~\eqref{eq:ch15:floor} of the low-weight codewords generated by
weight-2 input patterns, found by exhaustive search. The QPP interleaver has no weight-2 codeword
lighter than $d=28$; the random one has one of weight 14.}

Figure~\ref{fig:ch15_turbo_floor} shows the floor. In the waterfall region the two interleavers
are indistinguishable. Below about $10^{-4}$, however, the random interleaver's curve bends over and
falls by only a factor of two or three per half decibel: the error floor. The search for weight-2
input patterns finds a codeword of weight only 14 for the random interleaver, while the QPP
interleaver, whose algebraic structure spreads every pair of nearby positions, has none lighter than
28. The simulated floor of the random interleaver lies about an order of magnitude \emph{above} the
weight-2 estimate, a reminder that weight-3, weight-4 and higher input patterns, which the search
does not include, also produce low-weight codewords; the dashed curves are estimates of one
contribution, not bounds. Designing interleavers that remove all such patterns up to some weight is
the aim of the S-random, dithered relative prime (DRP), almost regular permutation (ARP) and QPP
designs.

\begin{worked}[How low is the floor?]
Suppose a rate-1/3 turbo code with $K=1024$ has two codewords of weight $d=14$, each produced by a
weight-2 input. Estimate the floor at $\ebno=2$\,dB.

$\ebno=10^{0.2}=1.585$, so the argument of the Q-function is
$\sqrt{2\cdot\frac13\cdot14\cdot1.585}=\sqrt{14.8}=3.85$, and $\Q(3.85)\approx5.9\times10^{-5}$.
Each such codeword causes two bit errors out of 1024 when it is mistaken for the transmitted one,
so~\eqref{eq:ch15:floor} gives
$P_b\approx2\times\frac{2}{1024}\times5.9\times10^{-5}\approx2.3\times10^{-7}$.
The floor falls by only about a factor of 4 per decibel (compare the waterfall, which falls by a
factor of 100 or more per decibel). A system that needs $10^{-10}$, such as an optical transport
link or a satellite broadcast of compressed video, either needs a code with no such codewords or an
outer code to clean up after the inner one---which is why DVB-S2 places a BCH code outside its LDPC
code, and why optical standards rarely use turbo codes.
\end{worked}

\subsection{EXIT charts}

Why is there a threshold at all, and where is it? The answer came in 2001 from Stephan ten Brink's
\term{extrinsic information transfer} (EXIT) chart, one of the most useful pictures in coding theory.
The idea is to characterise each SISO decoder by a single transfer function: how much information
about the bits comes \emph{out} of it, as a function of how much goes \emph{in}.

Measure the quality of a set of LLRs $L$ about bits $c$ by the mutual information $I(C;L)$, a number
between 0 (useless) and 1 (perfect). Ten Brink's key observation was that the a-priori LLRs a
decoder receives in the middle of the iterations are well modelled as consistent Gaussian variables,
$L_A=\frac{\sigma_A^2}{2}x+\mathcal N(0,\sigma_A^2)$, whose mutual information is a known function
of $\sigma_A$ alone:
\begin{equation}
I_A=J(\sigma_A)=1-\E\!\left[\log_2\!\big(1+e^{-L_A}\big)\right]_{x=+1}.
\label{eq:ch15:J}
\end{equation}
To measure a decoder's transfer curve $I_E=T(I_A)$ at a given channel SNR, feed it channel LLRs for a
long block and synthetic a-priori LLRs with a chosen $I_A$, run it once, and measure the mutual
information of its extrinsic output, conveniently estimated by the time average
$I_E\approx1-\frac1K\sum_k\log_2(1+e^{-x_kL_{E,k}})$. Figure~\ref{fig:ch15_exit}(a) shows such
curves for the LTE constituent decoder. Each starts above zero at $I_A=0$ (the decoder learns
something from the channel alone) and reaches 1 at $I_A=1$; a better channel lifts the whole curve.

In the turbo decoder, decoder 1's extrinsic output is decoder 2's a-priori input and vice versa. Plot
decoder 1's curve $I_{E1}=T_1(I_{A1})$ and decoder 2's curve with its axes swapped
($I_{A1}=I_{E2}=T_2(I_{A2})$ with $I_{A2}=I_{E1}$) on the same unit square
(Figure~\ref{fig:ch15_exit}b). The decoding trajectory is a staircase between the two curves:
vertical steps for decoder 1, horizontal steps for decoder 2. If the curves intersect, the staircase
gets stuck at the intersection and the decoder fails, however many iterations it runs. If there is an
open \term{tunnel} between them, the staircase climbs through it to $(1,1)$ and the block is decoded.
The threshold is the SNR at which the tunnel first opens, and the number of steps needed to get
through it---many when the tunnel is narrow---is the number of iterations. The measured trajectory of
a real decoder with a long random interleaver follows the predicted staircase closely.

\mdcfig{ch15_exit}{EXIT analysis of the rate-1/3 LTE turbo code. (a) Transfer curves of one
(13,15) RSC decoder (log-MAP) at five values of $\ebno$, each point averaged over 160\,000 bits with
Gaussian a-priori LLRs. (b) The EXIT chart at $+0.5$\,dB, where a tunnel is open, and at $-0.5$\,dB,
where the curves cross. The staircase is the trajectory measured inside a real iterative decoder
($K=10\,000$, random interleaver).}

The chart explains several empirical facts at once. Symmetric decoders' tunnel opens when a
transfer curve lies above the diagonal everywhere; for the LTE code this happens at about
$-0.1$\,dB, about 0.4\,dB from the rate-1/3 limit, and that is where the waterfall of a very long
code begins. Short blocks need more margin, because the Gaussian model and the independence
assumption hold only on average. A strong constituent code (many states) gives a curve that starts
low and ends steeply; a weak one gives a flatter curve---and it is the \emph{match} between the two
curves, not the strength of each, that determines the threshold. That insight turned code design
into curve fitting: choose component codes, degree distributions or bit mappings whose curves fit
together with a narrow but open tunnel.

\begin{keyidea}[The area theorem]
For the binary erasure channel, and approximately for other channels, the area under an EXIT curve
is fixed by the component code's rate (Ashikhmin, Kramer and ten Brink, 2004). A capacity-achieving
iterative system therefore needs the two curves to \emph{match}: any area between them is lost rate.
The narrower the tunnel, the closer to capacity, and the more iterations are needed to crawl
through it. Approaching capacity costs decoding effort, and the EXIT chart shows exactly how much.
\end{keyidea}

\subsection{Serial concatenation and the turbo principle}

In a \term{serially concatenated convolutional code} (SCCC; Benedetto, Divsalar, Montorsi and
Pollara, 1998) an outer code encodes the information, an interleaver permutes the result, and an
inner \emph{recursive} code encodes it again. The decoder iterates between a SISO decoder for the
inner code (which produces extrinsic information about the inner code's input, that is, the outer
code's coded bits) and one for the outer code. Serial concatenation typically has a slightly later
waterfall than parallel concatenation but a much lower error floor, because the outer code's
distance multiplies the interleaver gain: the floor falls as $K^{-\lfloor(d_o+1)/2\rfloor}$, where
$d_o$ is the outer code's free distance. The simplest members of the family are
\term{repeat--accumulate} (RA) codes (Divsalar, Jin and McEliece, 1998): repeat each bit $q$ times,
interleave, and pass the result through the accumulator $1/(1+D)$. RA codes are both turbo-like and
LDPC-like, and their irregular versions are the ancestors of the DVB-S2 and 5G LDPC codes.

The \term{turbo principle}---iterate between SISO modules that exchange extrinsic LLRs through
interleavers---applies whenever a system can be viewed as a concatenation. \term{Turbo equalization}
(Chapter~\ref{ch:equalization}) treats an ISI channel as an inner rate-1 code; iterative detection
and decoding of MIMO (Chapter~\ref{ch:mimo}) treats the spatial multiplexer as one; and
\term{bit-interleaved coded modulation} with iterative demapping (BICM-ID) treats the constellation
labelling as one. In each case the EXIT chart of the inner module against the outer decoder predicts
the gain.

\subsection{Turbo codes in practice}

\begin{table}[htbp]\centering\small
\caption{Turbo codes in standards. ``Duo-binary'' codes take two bits per trellis step and use
circular (tail-biting) termination.}
\label{tab:ch15:turbostd}
\begin{tabularx}{\linewidth}{@{}l l l >{\raggedright\arraybackslash}X@{}}
\toprule
System & Constituent code & Interleaver & Remarks\\
\midrule
UMTS/HSPA (TS 25.212) & 8-state RSC (13,15) & prime-based, $K\le5114$ & rate 1/3, 12 tail bits\\
LTE/LTE-A (TS 36.212) & 8-state RSC (13,15) & QPP, $40\le K\le6144$ & 188 sizes; circular-buffer rate matching\\
CCSDS 131.0-B (space) & 16-state & algebraic, $K$ = 1784--8920 & rates 1/2 to 1/6\\
DVB-RCS and RCS2 & duo-binary, 8/16 states & two-level permutation & satellite return links\\
IEEE 802.16e (WiMAX) & duo-binary, 8 states & ARP-type & convolutional turbo code option\\
HomePlug AV, IEEE 1901 & duo-binary, 8 states & tabulated & power-line communication\\
\bottomrule
\end{tabularx}
\end{table}

Table~\ref{tab:ch15:turbostd} lists the main deployments. Turbo codes were the first
capacity-approaching codes in mass production: every 3G and 4G handset contains a turbo decoder,
and for about fifteen years the LTE turbo decoder was one of the largest and most power-hungry
blocks in a modem's baseband. Its critical path is the BCJR recursion, which is inherently
sequential. Throughput is obtained by splitting each block into windows (typically 32 to 64 steps)
processed by parallel BCJR engines---made possible by the QPP interleaver's contention-free
property---with the boundary state metrics initialised from the previous iteration or from short
``training'' recursions. As LTE-Advanced pushed peak rates past 1\,Gb/s, the turbo decoder's area
and power, and the difficulty of parallelising it further, became a principal argument for a
different code in 5G.

\begin{inpractice}[Deep space: turbo codes near the limit]
The Consultative Committee for Space Data Systems added turbo codes to its telemetry recommendation
(CCSDS 131.0-B) in 1999, with information block lengths of 1784, 3568, 7136 and 8920 bits and rates
1/2, 1/3, 1/4 and 1/6. At the low rates they operate roughly 1\,dB from the capacity limit at
frame error rates of about $10^{-4}$, and more than 2\,dB better than the Reed--Solomon plus
convolutional concatenation that preceded them. For a spacecraft at Mars, 2\,dB is a 60\% increase
in data return from the same antenna and transmitter. NASA's Mars Reconnaissance Orbiter, Messenger
and New Horizons missions, among others, used CCSDS turbo codes. Deep space is an ideal home for
them: low rates, long blocks, modest data rates and an error-rate target above the floor.
\end{inpractice}

%======================================================================
\section{Low-density parity-check codes}
\label{sec:ch15:ldpc}
%======================================================================

\begin{history}[The code that was thirty-five years early]
Robert Gallager's 1960 MIT doctoral thesis, published as a short monograph by MIT Press in 1963,
defined \emph{low-density parity-check codes}: linear codes whose parity-check matrix has only a
few ones in each row and column. Gallager proved that such codes, chosen at random with fixed row
and column weights, have minimum distance growing linearly with the block length, and he described
iterative decoding algorithms---including one that passes probabilities between the bits and the
checks, which is belief propagation in all but name---that worked remarkably well in his
simulations. But the simulations had to run on the computers of 1960, and the concatenated codes of
Forney (1966) and the Viterbi algorithm (1967) soon offered more practical routes. LDPC codes were
largely forgotten. Michael Tanner's 1981 paper introduced the bipartite graph representation that now
bears his name and generalised it to codes built from component codes, but it too drew little
attention. In 1995--96, prompted by turbo codes, David MacKay and Radford Neal at Cambridge, and
independently Michael Sipser and Daniel Spielman at MIT, rediscovered sparse-graph codes. MacKay and
Neal's letter ``Near Shannon limit performance of low density parity check codes'' showed that
Gallager's codes, decoded by belief propagation, performed almost as well as turbo codes. Within five
years, Richardson, Shokrollahi and Urbanke had shown how to design \emph{irregular} LDPC codes that
beat turbo codes, and Chung, Forney, Richardson and Urbanke had simulated one within 0.04\,dB of the
Shannon limit. Gallager, asked about the rediscovery, remarked that he had not expected anyone to
need the codes for another few decades.
\end{history}

\subsection{Parity-check matrices and Tanner graphs}

A binary linear code of length $n$ and dimension $k$ is the null space of an $m\times n$
parity-check matrix $\vect H$ ($m\ge n-k$): $\vect c$ is a codeword if and only if $\vect H\vect c^T=\vect 0$
over GF(2) (Chapter~\ref{ch:classiccodes}). An \term{LDPC code} is one whose $\vect H$ is sparse: the
number of ones grows linearly with $n$, not with $n^2$. A $(d_v,d_c)$-\emph{regular} code has
exactly $d_v$ ones in every column and $d_c$ in every row; Gallager's favourite was $(3,6)$, which
has design rate
\begin{equation}
R\ge1-\frac{m}{n}=1-\frac{d_v}{d_c}=\frac12.
\end{equation}
(The inequality allows for redundant rows.) For $n=2304$, the $(3,6)$ code's $\vect H$ has
$1152\times2304\approx2.65$ million entries but only 6912 ones.

The \term{Tanner graph} of $\vect H$ is a bipartite graph with a \term{variable node} for each
code bit (each column), a \term{check node} for each parity check (each row), and an edge wherever
$H_{ij}=1$ (Figure~\ref{fig:ch15_tanner}). The code is the set of assignments to the variable nodes
that make every check node's neighbourhood have even parity. The number of edges, $E=nd_v=md_c$,
measures the decoding complexity per iteration. A \term{cycle} is a closed path in the graph; its
length is even and at least 4, and the length of the shortest cycle is the \term{girth}.

\begin{figure}[htbp]\centering
\begin{tikzpicture}[x=1.05cm,y=1cm,
  vn/.style={circle,draw=navy,fill=navy!10,thick,minimum size=6mm,font=\sffamily\scriptsize},
  cn/.style={rectangle,draw=accent,fill=accent!10,thick,minimum size=5.5mm,font=\sffamily\scriptsize}]
\begin{scope}[shift={(-5.2,1.1)}]
\node[font=\sffamily\small,navy] at (1.4,1.1) {$\vect H$};
\node[font=\small] at (1.4,0) {$\begin{bmatrix}1&1&0&1&1&0&0\\ 1&0&1&1&0&1&0\\ 0&1&1&1&0&0&1\end{bmatrix}$};
\node[font=\sffamily\scriptsize] at (1.4,-1.1) {columns $c_0,\dots,c_6$};
\end{scope}
\foreach \i in {0,...,6}{ \node[vn] (v\i) at (\i*1.05,0) {$c_{\i}$}; }
\foreach \j/\x in {0/1.05,1/3.15,2/5.25}{ \node[cn] (c\j) at (\x,2.4) {$\oplus$}; }
\foreach \j/\list in {0/{0,1,3,4},1/{0,2,3,5},2/{1,2,3,6}}{
  \foreach \i in \list { \draw[navy!60,thick] (c\j.south) -- (v\i.north); } }
\draw[accent,line width=2pt,opacity=0.55] (c0.south) -- (v0.north) -- (c1.south) -- (v3.north) -- cycle;
\node[font=\sffamily\scriptsize,accent,anchor=west] at (6.9,1.6) {a 4-cycle:};
\node[font=\sffamily\scriptsize,accent,anchor=west] at (6.9,1.25) {$c_0,c_3$ share};
\node[font=\sffamily\scriptsize,accent,anchor=west] at (6.9,0.9) {two checks};
\node[font=\sffamily\scriptsize,navy,anchor=east] at (-0.4,0) {variable nodes};
\node[font=\sffamily\scriptsize,accent,anchor=east] at (0.6,2.4) {check nodes};
\end{tikzpicture}
\caption{The Tanner graph of a parity-check matrix (here the $(7,4)$ Hamming code of
Chapter~\ref{ch:classiccodes}, which is far too dense to be a good LDPC code but small enough to
draw). Each column is a variable node (circle), each row a check node (square), each 1 an edge. The
highlighted 4-cycle arises because columns 0 and 3 overlap in two rows; such short cycles degrade
belief propagation, and good LDPC constructions avoid them.}
\label{fig:ch15_tanner}
\end{figure}

\subsection{Degree distributions and irregular codes}

In an \term{irregular} LDPC code the node degrees vary. The standard description counts
\emph{edges}: $\lambda_i$ is the fraction of edges attached to variable nodes of degree $i$, and
$\rho_j$ the fraction attached to check nodes of degree $j$, collected in the polynomials
\begin{equation}
\lambda(x)=\sum_i\lambda_i x^{i-1},\qquad\rho(x)=\sum_j\rho_jx^{j-1}.
\label{eq:ch15:degpoly}
\end{equation}
(The exponent $i-1$ is the number of \emph{other} edges at a node, which is what matters in message
passing.) The $(3,6)$ code has $\lambda(x)=x^2$ and $\rho(x)=x^5$. The design rate is
\begin{equation}
R=1-\frac{\int_0^1\rho(x)\,dx}{\int_0^1\lambda(x)\,dx}=1-\frac{\sum_j\rho_j/j}{\sum_i\lambda_i/i}.
\label{eq:ch15:designrate}
\end{equation}
Why be irregular? High-degree variable nodes collect many opinions and become reliable quickly,
then help their neighbours; low-degree nodes are cheap in rate. A well-chosen mix creates a
\emph{wave} of reliability that sweeps from the high-degree nodes through the rest of the graph.
Optimised irregular codes have thresholds within hundredths of a decibel of capacity, while regular
codes stall about 1\,dB away (Section~\ref{sec:ch15:de}).

\begin{worked}[Rate of an irregular ensemble]
An ensemble has $\lambda(x)=0.3x+0.3x^2+0.4x^7$ (edges on degree-2, 3 and 8 variable nodes) and
$\rho(x)=x^6$ (all checks of degree 7). What is its design rate, and what fraction of variable nodes
have degree 2?

$\sum_i\lambda_i/i=0.3/2+0.3/3+0.4/8=0.15+0.10+0.05=0.30$ and $\sum_j\rho_j/j=1/7=0.1429$, so
$R=1-0.1429/0.30=0.524$. The fraction of \emph{nodes} of degree $i$ is proportional to
$\lambda_i/i$: degree-2 nodes are $0.15/0.30=50\%$ of the variable nodes even though they carry only
30\% of the edges. This conversion between edge and node perspectives is a frequent source of
errors.
\end{worked}

\subsection{Belief propagation}

Belief propagation (BP), also called the \term{sum-product algorithm}, computes on the Tanner graph
exactly what the BCJR algorithm computes on a trellis: the marginal APP of each bit, under the
assumption that the graph has no cycles. Messages are LLRs passed along edges, in two kinds of
update.

\paragraph{Variable-node update.} A variable node $v$ with channel LLR $L_v$ receives messages
$r_{c\to v}$ from its check neighbours $\mathcal N(v)$. Each incoming message is independent
evidence about the same bit (on a tree), so the node's belief is their sum, and the message it
sends back to check $c$ excludes what $c$ told it:
\begin{equation}
q_{v\to c}=L_v+\sum_{c'\in\mathcal N(v)\setminus c}r_{c'\to v},\qquad
L^{\text{APP}}_v=L_v+\sum_{c'\in\mathcal N(v)}r_{c'\to v}.
\label{eq:ch15:vn}
\end{equation}

\paragraph{Check-node update.} A check node $c$ says that the XOR of its neighbours' bits is zero,
so the bit at $v$ equals the XOR of all the others. By the box-plus rule~\eqref{eq:ch15:boxplus},
\begin{equation}
r_{c\to v}=\mathop{\boxplus}_{v'\in\mathcal N(c)\setminus v}q_{v'\to c}
=2\tanh^{-1}\!\prod_{v'\in\mathcal N(c)\setminus v}\tanh\frac{q_{v'\to c}}{2},
\label{eq:ch15:tanh}
\end{equation}
the \term{tanh rule}. It is often written in a form that separates signs and magnitudes, using the
function $\phi(x)=-\ln\tanh(x/2)$, which is its own inverse for $x>0$
(Figure~\ref{fig:ch15_checknode}a):
\begin{equation}
r_{c\to v}=\Big(\prod_{v'}\operatorname{sign}q_{v'\to c}\Big)\;
\phi\Big(\sum_{v'}\phi\big(|q_{v'\to c}|\big)\Big).
\label{eq:ch15:phi}
\end{equation}
The product of tanh becomes a sum in the $\phi$ domain, so hardware can implement the magnitude path
with two lookup tables and an adder tree.

\paragraph{The algorithm.} Initialise all check-to-variable messages to zero; then iterate: update
all variable nodes, update all check nodes, and after each iteration make hard decisions from
$L^{\text{APP}}$. Stop when the decisions satisfy $\vect H\hat{\vect c}^T=\vect 0$ (a free and very
reliable stopping rule) or after a maximum number of iterations, typically 10 to 50. The cost per
iteration is a few operations per edge.

If the Tanner graph were a tree, BP would compute the exact APPs. Real codes have cycles, so after
a few iterations the messages are no longer independent and BP becomes an approximation. For graphs
whose cycles are long, the approximation is excellent: the neighbourhood of a node looks like a tree
for $\lfloor(g-2)/4\rfloor$ iterations if the girth is $g$, and random-like dependencies thereafter
do surprisingly little harm. Short cycles, particularly 4-cycles, cause a node's own message to
return to it after two iterations, and constructions such as \term{progressive edge growth}
(PEG; Hu, Eleftheriou and Arnold, 2005; implemented in \texttt{commlib.LDPCCode}) and the shift
selection of quasi-cyclic codes are designed to avoid them.

\begin{keyidea}[Turbo decoding is belief propagation]
In 1998 McEliece, MacKay and Cheng, and independently Kschischang and Frey, showed that the turbo
decoder is an instance of Judea Pearl's belief propagation algorithm (1988) on a graph that contains
the two trellises and the interleaver, and that Gallager's LDPC decoder is the same algorithm on a
simpler graph. The factor-graph framework of Kschischang, Frey and Loeliger (2001) unified BCJR,
Viterbi, turbo decoding, LDPC decoding, the Kalman filter and the FFT as message passing on graphs.
Once you see the graph, the decoder follows.
\end{keyidea}

\subsection{Min-sum and its corrections}

The tanh rule needs transcendental functions; hardware designers wanted comparators. Applying the
approximation~\eqref{eq:ch15:minapprox} to every pair gives the \term{min-sum} check update,
\begin{equation}
r_{c\to v}\approx\Big(\prod_{v'\ne v}\operatorname{sign}q_{v'\to c}\Big)\min_{v'\ne v}|q_{v'\to c}|.
\label{eq:ch15:minsum}
\end{equation}
A check node therefore needs only the signs, the smallest magnitude, the second smallest and the
position of the smallest: every outgoing message is the smallest magnitude of the \emph{others},
which is the overall minimum except for the edge that supplied it, which receives the second minimum.
Min-sum needs no knowledge of the noise variance (scaling all inputs scales all messages), which is
convenient, but its outputs are systematically too large, as Figure~\ref{fig:ch15_checknode}(b) shows:
the exact output is always smaller in magnitude than the minimum. Two one-parameter corrections
recover most of the loss:
\begin{itemize}
\item \term{normalised min-sum} (NMS) multiplies the magnitude by $\alpha\approx0.7$--$0.85$;
\item \term{offset min-sum} (OMS) subtracts $\beta\approx0.3$--$0.5$ and clips at zero.
\end{itemize}
With tuned parameters both come within about 0.1 to 0.2\,dB of sum-product for most practical codes, and almost every LDPC decoder
in silicon today, in phones, Wi-Fi chips, SSD controllers and optical DSPs, uses one of them with
messages quantised to 4 to 6 bits.

\mdcfig{ch15_checknode}{(a) The function $\phi(x)=-\ln\tanh(x/2)$ used in the check-node update;
it is self-inverse, very large for small arguments and very small for large ones, so the sum in
\eqref{eq:ch15:phi} is dominated by the least reliable input. (b) Outputs of a degree-6 check node
(five inputs drawn from $\mathcal N(3.5,2.5^2)$): the exact tanh rule against plain min-sum. Min-sum
is always over-confident; normalisation and offset pull its output back towards the exact value.}

\mdcfig{ch15_ldpc_decoders}{(a) Frame (solid) and bit (dashed) error rates of a $(3,6)$-regular
quasi-cyclic LDPC code with $n=2304$ (girth-8 circulant shifts, $Z=384$) under four check-node rules,
50 flooding iterations, all-zero codeword (valid because the code and decoders are symmetric), at
least 60 frame errors or 1500 frames per point. (b) Mean number of iterations to satisfy all checks
with normalised min-sum, for the flooding and the layered schedule (500 frames per point).}

Figure~\ref{fig:ch15_ldpc_decoders}(a) compares the four rules on a regular $(3,6)$ code. Plain
min-sum loses about half a decibel at a frame error rate of $10^{-2}$, normalised and offset min-sum lose about 0.2\,dB (with these untuned parameters), and all four
curves sit to the right of the ensemble's BP threshold of 1.11\,dB, which in turn is about
0.9\,dB from the rate-1/2 limit. That last gap is the price of regularity, and irregular designs
remove most of it.

\subsection{Schedules: flooding and layered decoding}

The algorithm above updates all variable nodes, then all check nodes, a \term{flooding} schedule.
The \term{layered} (or serial, or shuffled) schedule instead processes the check nodes in groups,
called layers, and updates the posterior LLRs of the variables involved immediately after each
group, so that later layers in the same iteration already use the new information. Layered decoding
of a check $c$ is
\begin{equation}
q_{v\to c}=L^{\text{APP}}_v-r^{\text{old}}_{c\to v},\qquad
r^{\text{new}}_{c\to v}=\mathop{\boxplus}_{v'\ne v}q_{v'\to c},\qquad
L^{\text{APP}}_v\leftarrow q_{v\to c}+r^{\text{new}}_{c\to v},
\end{equation}
and needs no separate storage for the variable-to-check messages. It converges in roughly half as
many iterations as flooding (Figure~\ref{fig:ch15_ldpc_decoders}b), which halves the latency and
the energy per decoded bit, and it is the standard architecture for quasi-cyclic codes, whose block
rows form natural layers in which each variable appears at most once.

\subsection{Decoding on the erasure channel: peeling}

On the \term{binary erasure channel} (BEC) each bit arrives either correctly or as an erasure ``?'',
and BP reduces to a simple combinatorial algorithm. A check node with exactly one erased neighbour
determines it: the erased bit is the XOR of the others. Resolve it, remove it from the graph, and
repeat. This \term{peeling decoder} either recovers every erasure or stops at a \term{stopping set},
a set of erased variables every one of whose neighbouring checks sees at least two of them. Peeling
is exactly BP on the BEC and costs one operation per edge in total. It is the decoder of the fountain
and Raptor codes used for file delivery in 3GPP multimedia broadcast and DVB (Luby and Shokrollahi),
and the same successive-resolution idea drives the contention-resolution random access of
Chapter~\ref{ch:satellite}.

%======================================================================
\section{Density evolution: computing the threshold}
\label{sec:ch15:de}
%======================================================================

How good can an LDPC code with given degree distributions be, in the limit of long blocks? Richardson
and Urbanke answered in 2001 with \term{density evolution} (DE). Their concentration theorem says
that the performance of a randomly chosen code from an ensemble concentrates around the ensemble
average as $n$ grows, and that the average behaves like BP on a cycle-free graph. On a tree, the
messages at each iteration are independent, so it suffices to track their \emph{probability
distribution}. There is a threshold channel parameter below which the error probability tends to
zero with the iterations and above which it does not: the \term{BP threshold}.

\subsection{The erasure channel: a closed form}

On the BEC with erasure probability $\varepsilon$, a message is either correct or erased, so its
``density'' is a single number: the probability $x_\ell$ that a variable-to-check message is erased
at iteration $\ell$. A check-to-variable message is erased unless all $d_c-1$ other incoming messages
are known, with probability $1-(1-x)^{d_c-1}$; averaged over the check-degree distribution this is
$1-\rho(1-x)$. A variable-to-check message is erased only if the channel erased the bit \emph{and}
all $d_v-1$ other incoming check messages are erased. Hence
\begin{equation}
x_{\ell+1}=\varepsilon\,\lambda\big(1-\rho(1-x_\ell)\big),\qquad x_0=\varepsilon.
\label{eq:ch15:debec}
\end{equation}
Decoding succeeds if $x_\ell\to0$, which happens if and only if
$\varepsilon\lambda(1-\rho(1-x))<x$ for all $x\in(0,\varepsilon]$. The threshold is therefore
\begin{equation}
\varepsilon^*=\min_{x\in(0,1]}\frac{x}{\lambda\big(1-\rho(1-x)\big)}.
\label{eq:ch15:becthr}
\end{equation}
Figure~\ref{fig:ch15_de_bec}(a) shows the recursion as a staircase for the $(3,6)$ ensemble. Below the
threshold the curve $\varepsilon\lambda(1-\rho(1-x))$ lies under the diagonal and the staircase runs
down to zero; above it, the curve touches the diagonal and the staircase stops at a nonzero fixed
point, where a fraction of the bits stays erased forever. The plot is an EXIT chart in disguise.

\mdcfig{ch15_de_bec}{(a) Density evolution for the $(3,6)$ ensemble on the BEC: the map
$x\mapsto\varepsilon\lambda(1-\rho(1-x))$ below, at and above the threshold $\varepsilon^*=0.4294$,
with the decoding staircases. (b) Residual erasure rate of the peeling decoder on randomly lifted
$(3,6)$ codes of increasing length (400 to 4 frames per point). As $n$ grows the curves steepen
towards the threshold predicted by density evolution; capacity would allow $\varepsilon=0.5$.}

\begin{worked}[The $(3,6)$ threshold on the BEC]
For $\lambda(x)=x^2$, $\rho(x)=x^5$, the fixed-point condition is
$\varepsilon\big(1-(1-x)^5\big)^2=x$. Find $\varepsilon^*$.

Write $f(x)=x/(1-(1-x)^5)^2$ and minimise it over $x$. At $x=0.2$: $(0.8)^5=0.328$, so
$f=0.2/0.672^2=0.443$. At $x=0.26$: $(0.74)^5=0.222$, $f=0.26/0.778^2=0.4296$. At $x=0.3$:
$(0.7)^5=0.168$, $f=0.3/0.832^2=0.433$. The minimum is near $x=0.26$ with value
$\varepsilon^*\approx0.4294$ (a finer search confirms 0.42944). A rate-1/2 code could in principle
tolerate $\varepsilon=0.5$, so the $(3,6)$ ensemble under BP achieves 86\% of capacity on the BEC. The
$(4,8)$ ensemble, with the same rate, does worse ($\varepsilon^*\approx0.383$): more edges make the
graph more constrained but make the early iterations harder. Irregular ensembles such as
$\lambda(x)=0.5x+0.5x^{2}$ paired with suitable check degrees can come arbitrarily close to 0.5.
\end{worked}

Figure~\ref{fig:ch15_de_bec}(b) shows that real finite codes approach the DE prediction. The error
curves of the peeling decoder on lifted $(3,6)$ codes steepen as $n$ grows, and for $n$ in the
hundreds of thousands the transition is within a hundredth of the threshold. The slope at finite
length follows a scaling law in $\sqrt n\,(\varepsilon-\varepsilon^*)$, the LDPC analogue of the
channel dispersion of Chapter~\ref{ch:infotheory}.

\subsection{The AWGN channel}

On the BI-AWGN channel the messages are real-valued LLRs and DE must track entire probability
densities. With the all-zero codeword assumed (justified by symmetry), the channel LLR has density
$\mathcal N(2/\sigma^2,4/\sigma^2)$. The variable-node update adds independent LLRs, so densities
\emph{convolve}; the check-node update adds independent variables in the $\phi$ domain, so densities
convolve after a change of variable. Richardson and Urbanke implemented these convolutions
numerically with FFTs on quantised densities. A simpler alternative, used in
Figure~\ref{fig:ch15_de_awgn}, is \emph{Monte Carlo DE}: represent each density by a large
population of samples and apply the update rules to randomly drawn members.

\mdcfig{ch15_de_awgn}{Density evolution for the $(3,6)$ ensemble on the BI-AWGN channel, by
population dynamics with 200\,000 samples. (a) The density of variable-to-check LLRs at
$\sigma=0.82$ drifts to the right and widens, iteration by iteration, as it must for consistent
densities (variance twice the mean). (b) Bit error probability versus iteration: below the threshold
$\sigma^*\approx0.881$ ($\ebno\approx1.10$\,dB) the error probability eventually collapses, more
slowly the closer the channel is to the threshold; above it, it stalls.}

The result for the $(3,6)$ ensemble is $\sigma^*\approx0.881$, that is, $\ebno^*\approx1.10$\,dB,
compared with the BPSK capacity limit of 0.19\,dB. Figure~\ref{fig:ch15_de_awgn}(b) shows the
characteristic behaviour near the threshold: at $\sigma=0.87$ the decoder spends thirty iterations
in a ``bottleneck'' where the messages improve only slowly---the DE curve passing close to the
diagonal---and then breaks through.

\paragraph{The Gaussian approximation.} Chung, Richardson and Urbanke (2001) observed that the
densities stay close to consistent Gaussians, $\mathcal N(\mu,2\mu)$, so that one number, the mean,
suffices. At a degree-$d_v$ variable node means add; at a degree-$d_c$ check node the tanh rule
becomes, in expectation,
\begin{equation}
\mu_c=\phi_G^{-1}\!\Big(1-\big[1-\phi_G(\mu_v)\big]^{d_c-1}\Big),\qquad
\phi_G(\mu)=1-\E\big[\tanh(L/2)\big]_{L\sim\mathcal N(\mu,2\mu)},
\label{eq:ch15:ga}
\end{equation}
with simple fitted approximations for $\phi_G$. The \term{Gaussian approximation} (GA) predicts
thresholds within about 0.1\,dB and turns degree-distribution optimisation into a fast linear program
(in $\lambda$, for fixed $\rho$). It is also, as we will see, the standard way to construct polar
codes. EXIT charts are the same idea expressed in mutual information instead of mean.

\begin{inpractice}[How close is close?]
Richardson, Shokrollahi and Urbanke's optimised rate-1/2 ensembles with maximum variable degree 100
have BI-AWGN thresholds of $\sigma^*\approx0.97$, about 0.06\,dB from capacity; the ensemble of Chung
\emph{et al.} with maximum degree 8000 has a threshold within 0.0045\,dB. A code with $10^7$ bits
drawn from it reached $10^{-6}$ within 0.04\,dB of the limit. Such codes are laboratory records: their
very high degrees make them impractical, and their finite-length behaviour at moderate lengths is
worse than that of more modest designs. Practical codes, with maximum variable degree around 10 to
30 and structured graphs, have thresholds about 0.2 to 0.5\,dB from capacity and, at lengths of a
few thousand bits, operate about 1\,dB from the limit at a frame error rate of $10^{-2}$.
\end{inpractice}

%======================================================================
\section{Structured LDPC codes and their standards}
\label{sec:ch15:qc}
%======================================================================

\subsection{Encoding}

A random sparse $\vect H$ has a dense generator matrix, so encoding by $\vect c=\vect u\vect G$ costs
$O(n^2)$ operations and storage---about $1.3$ million XORs and bits of memory for $n=2304$, and 5
billion for $n=64\,800$. (The encoder in \texttt{commlib.LDPCCode} does exactly this, which is fine
for teaching.) Practical codes are designed to be encodable directly from $\vect H$. The most common
structure splits $\vect H=[\vect H_u\;\vect H_p]$ into an information part and a parity part, with
$\vect H_p$ lower triangular or \emph{dual diagonal} (ones on the diagonal and just below it):
\begin{equation}
\vect H_p=\begin{bmatrix}1&&&\\1&1&&\\&1&1&\\&&\ddots&\ddots\end{bmatrix}.
\end{equation}
The parity bits are then computed one at a time by back-substitution, $p_i=p_{i-1}\oplus
(\vect H_u\vect u^T)_i$: an \emph{accumulator}. This is the irregular repeat--accumulate structure of
DVB-S2. Richardson and Urbanke showed in 2001 that any LDPC matrix can be brought to an
approximately triangular form with encoding complexity essentially linear in $n$.

\subsection{Protographs and quasi-cyclic lifting}

Modern LDPC codes are not random. They are built in two steps. First, design a small
\term{protograph} (or base graph): a Tanner graph with a few tens of nodes, possibly with parallel
edges, whose structure determines the threshold---DE and EXIT analysis extend naturally to
protographs (Thorpe, 2003). Second, \emph{lift} it: make $Z$ copies of the protograph and permute the
endpoints of each edge type among the copies. If each permutation is a cyclic shift, the result is a
\term{quasi-cyclic} (QC) LDPC code, whose parity-check matrix is a $m_b\times n_b$ array of
$Z\times Z$ blocks, each either zero or a cyclically shifted identity $\vect P^{s}$
(Figure~\ref{fig:ch15_lifting}). The whole code is described by the \term{base matrix} of shift values
$s_{ij}$ (with $-1$ for a zero block).

\begin{figure}[htbp]\centering
\begin{tikzpicture}[
  vn/.style={circle,draw=navy,fill=navy!10,thick,minimum size=5mm,inner sep=0pt,font=\sffamily\tiny},
  cn/.style={rectangle,draw=accent,fill=accent!10,thick,minimum size=4.5mm,inner sep=0pt,font=\sffamily\tiny}]
% protograph
\begin{scope}
\foreach \j/\x in {0/0.0,1/0.9,2/1.8,3/2.7}{ \node[vn] (p\j) at (\x,0) {$v_\j$}; }
\node[cn] (q0) at (0.9,1.6) {$c_0$}; \node[cn] (q1) at (1.8,1.6) {$c_1$};
\draw[navy!60,thick] (q0.south) -- (p0.north); \draw[navy!60,thick] (q0.south) -- (p1.north); \draw[navy!60,thick] (q0.south) -- (p2.north);
\draw[navy!60,thick] (q1.south) -- (p0.north); \draw[navy!60,thick] (q1.south) -- (p2.north); \draw[navy!60,thick] (q1.south) -- (p3.north);
\node[font=\sffamily\scriptsize,navy,align=center] at (1.35,-0.8) {protograph\\(2 checks, 4 variables)};
\node[font=\sffamily\scriptsize,align=center] at (1.35,2.45) {base matrix\\$\begin{bmatrix}0&1&3&-1\\2&-1&0&1\end{bmatrix}$};
\end{scope}
\draw[-{Stealth[length=3mm]},very thick,navy] (3.4,0.8) -- node[above,font=\sffamily\scriptsize]{lift, $Z=4$} (4.8,0.8);
% lifted H 8x16
\begin{scope}[shift={(5.2,2.35)},x=0.21cm,y=0.21cm]
\draw[gray!50] (0,0) rectangle (16,-8);
\foreach \k in {4,8,12}{ \draw[gray!40] (\k,0) -- (\k,-8); }
\draw[gray!40] (0,-4) -- (16,-4);
\foreach \bi/\bj/\s in {0/0/0,0/1/1,0/2/3,1/0/2,1/2/0,1/3/1}{
  \foreach \r in {0,...,3}{
    \pgfmathtruncatemacro{\c}{mod(\r+\s,4)}
    \fill[navy] (\bj*4+\c+0.12,-\bi*4-\r-0.12) rectangle (\bj*4+\c+0.88,-\bi*4-\r-0.88);
  } }
\node[font=\sffamily\scriptsize,anchor=west] at (16.6,-1.8) {$\vect P^{0}\;\vect P^{1}\;\vect P^{3}\;\vect 0$};
\node[font=\sffamily\scriptsize,anchor=west] at (16.6,-5.8) {$\vect P^{2}\;\vect 0\;\;\vect P^{0}\;\vect P^{1}$};
\node[font=\sffamily\scriptsize,align=center] at (8,-10.3) {lifted $\vect H$ ($8\times16$): each edge of the\\protograph becomes a $4\times4$ circulant};
\end{scope}
\end{tikzpicture}
\caption{Protograph lifting. A small base graph (left) is copied $Z$ times and each edge type is
replaced by a cyclic permutation among the copies; the shift values form the base matrix. The
lifted code keeps the degree structure and hence (for large $Z$) the threshold of the protograph,
while the shifts are chosen to avoid short cycles. A 4-cycle through base-matrix entries
$s_{i_1j_1},s_{i_1j_2},s_{i_2j_2},s_{i_2j_1}$ exists in the lifted graph exactly when
$s_{i_1j_1}-s_{i_1j_2}+s_{i_2j_2}-s_{i_2j_1}\equiv0\pmod Z$.}
\label{fig:ch15_lifting}
\end{figure}

QC structure is what made LDPC codes an industrial technology. The whole $\vect H$ is stored as a few
hundred small integers. The decoder processes a block row at a time with $Z$ identical processing
units, and the routing between memory and processors is a set of barrel shifters rather than an
arbitrary interconnect. The same base matrix serves many block lengths by changing $Z$, and the
shifts can be chosen by computer search to maximise the girth. The $n=2304$ code of
Figure~\ref{fig:ch15_ldpc_decoders} is an all-ones $3\times6$ protograph lifted with $Z=384$ and
shifts chosen at random until no 4- or 6-cycles remained.

\subsection{5G NR: two base graphs for everything}

The 5G New Radio data channels (PDSCH and PUSCH) use QC-LDPC codes defined in 3GPP TS 38.212 by just
two base graphs (Figure~\ref{fig:ch15_nr_bg}):
\begin{itemize}
\item \textbf{Base graph 1} (BG1) has $46\times68$ entries, 22 systematic columns, and is used for
large blocks and high rates. The largest code block has $K=22\times384=8448$ information bits.
\item \textbf{Base graph 2} (BG2) has $42\times52$ entries and 10 systematic columns, and is used for
small blocks and low rates, up to $K=10\times384=3840$.
\end{itemize}
Both have the same architecture. The top-left \emph{core}---4 rows, the systematic columns and 4 parity
columns with a dual-diagonal-like structure---is a high-rate code on its own (rate $22/24\approx0.92$
for BG1 after puncturing). Each further row adds one \emph{extension} parity bit that is a single
parity check on a few earlier bits, attached through an identity column (degree 1). The code is
therefore \emph{rate compatible} by construction: transmitting more of the extension parity bits
lowers the rate smoothly down to $22/66=1/3$ for BG1 and $10/50=1/5$ for BG2, exactly what hybrid ARQ
with incremental redundancy needs (Section~\ref{sec:ch15:harq}). The first two systematic columns
have very high degree and are always \emph{punctured} (never transmitted); these high-degree punctured
nodes act as hubs that improve the threshold, a trick borrowed from the protograph designs of
JPL's AR4JA codes.

\begin{figure}[htbp]\centering
\begin{tikzpicture}[x=0.1cm,y=0.1cm]
% BG1: 68 columns x 46 rows
\fill[navy!25] (0,0) rectangle (22,-4);
\fill[accent!30] (22,0) rectangle (26,-4);
\fill[navy!12] (0,-4) rectangle (26,-46);
\draw[practc,thick] (26,-4) -- (68,-46);
\fill[navy!45] (0,0) rectangle (2,-46);
\draw[navy,thick] (0,0) rectangle (68,-46);
\draw[navy!60] (22,0) -- (22,-46); \draw[navy!60] (26,0) -- (26,-46); \draw[navy!60] (0,-4) -- (68,-4);
\node at (12,3) {{\sffamily\scriptsize 22 systematic}};
\node at (24,6.5) {{\sffamily\scriptsize 4}};
\node at (47,3) {{\sffamily\scriptsize 42 extension parity}};
\node[anchor=east] at (-1,-2) {{\sffamily\scriptsize core: 4}};
\node[anchor=east] at (-1,-25) {{\sffamily\scriptsize 42}};
\node[align=center] at (12,-2.1) {{\sffamily\scriptsize $\vect A$}};
\node at (24,-2.1) {{\sffamily\scriptsize $\vect B$}};
\node at (13,-25) {{\sffamily\scriptsize $\vect C$ (sparse)}};
\node at (55,-22) {{\sffamily\scriptsize $\vect I$ (identity)}};
\node at (47,-2.1) {{\sffamily\scriptsize $\vect 0$}};
\node at (34,-52) {\sffamily\small\bfseries\color{navy} BG1: $46\times68$, $K_b=22$};
\draw[-{Stealth[length=1.5mm]},navy!70] (1,-49) -- (1,-46.5);
\node[anchor=west] at (-2,-51.5) {{\sffamily\scriptsize punctured}};
% BG2
\begin{scope}[shift={(85,0)}]
\fill[navy!25] (0,0) rectangle (10,-4);
\fill[accent!30] (10,0) rectangle (14,-4);
\fill[navy!12] (0,-4) rectangle (14,-42);
\draw[practc,thick] (14,-4) -- (52,-42);
\fill[navy!45] (0,0) rectangle (2,-42);
\draw[navy,thick] (0,0) rectangle (52,-42);
\draw[navy!60] (10,0) -- (10,-42); \draw[navy!60] (14,0) -- (14,-42); \draw[navy!60] (0,-4) -- (52,-4);
\node at (5,3) {{\sffamily\scriptsize 10}};
\node at (12,6.5) {{\sffamily\scriptsize 4}};
\node at (33,3) {{\sffamily\scriptsize 38 extension parity}};
\node[anchor=west] at (53,-2) {{\sffamily\scriptsize 4}};
\node[anchor=west] at (53,-23) {{\sffamily\scriptsize 38}};
\node at (26,-48) {\sffamily\small\bfseries\color{navy} BG2: $42\times52$, $K_b=10$};
\end{scope}
\end{tikzpicture}
\caption{Structure of the 5G NR LDPC base graphs (TS 38.212, Tables 5.3.2-2 and 5.3.2-3). The core
($\vect A$, $\vect B$) is a high-rate code with a dual-diagonal-like parity part $\vect B$ that allows
linear-time encoding; each extension row adds one degree-1 parity bit (the identity part), so
cutting the matrix at any row gives a lower- or higher-rate code. The first two systematic columns
(dark) have high degree and are never transmitted. Each nonzero entry is lifted to a $Z\times Z$
circulant, $Z\in\{2,\dots,384\}$.}
\label{fig:ch15_nr_bg}
\end{figure}

\paragraph{Lifting sizes.} The lifting size is $Z=a\times2^j$ with
$a\in\{2,3,5,7,9,11,13,15\}$ and $Z\le384$, giving 51 values from 2 to 384. They are grouped into eight
sets by $a$, each with its own table of shift coefficients $V_{ij}$, and the shift actually used is
$s_{ij}=V_{ij}\bmod Z$. One pair of base graphs and eight shift tables thus define several thousand
codes with fine granularity in length and rate, all decodable by the same hardware.

\paragraph{Choosing the base graph.} TS 38.212 selects BG2 when the transport block is small or the
rate low: if $A\le292$, or if $A\le3824$ and $R\le0.67$, or if $R\le0.25$; otherwise BG1
(Figure~\ref{fig:ch15_nr_basegraph}). A transport block CRC (24 bits if $A>3824$, otherwise 16) is
appended, and if the result exceeds the maximum code-block size (8448 for BG1, 3840 for BG2) it is
segmented into code blocks, each with its own 24-bit CRC.

\mdcfig[0.8]{ch15_nr_basegraph}{Base-graph selection for the 5G NR data channels as a function of
transport-block size $A$ and code rate $R$ (TS 38.212, Section 7.2.2). BG2 covers small blocks and
low rates, BG1 large blocks at medium and high rates.}

\paragraph{Rate matching.} The encoded bits (excluding the $2Z$ punctured systematic bits) are
written into a \emph{circular buffer}. Each transmission reads $E$ consecutive bits from it (skipping
filler bits), starting at one of four positions set by the \term{redundancy version}: for BG1 the
starting points are at approximately $0$, $17/66$, $33/66$ and $56/66$ of the buffer, for BG2 at $0$,
$13/50$, $25/50$ and $43/50$ (rounded to multiples of $Z$). RV0 starts at the systematic bits; RV2 and
RV3 start deep in the parity. If $E$ is smaller than the buffer the code is punctured; if larger, the
buffer wraps round and bits are repeated. The bits are then interleaved across the bit positions of
the QAM symbols (a row--column interleaver with as many rows as bits per symbol) and mapped.

\begin{worked}[From a transport block to a 5G LDPC codeword]
A PDSCH transport block has $A=8000$ bits and is scheduled with target code rate $R=0.5$. Follow it
through TS 38.212.

Since $A>3824$ and $R>0.25$ (and $R\le0.67$ does not matter because $A>3824$), \textbf{BG1} is used.
The transport block CRC is 24 bits (as $A>3824$), giving $B=8024$ bits. This is below 8448, so there
is a single code block and no code-block CRC. For BG1, $K_b=22$, and the lifting size is the smallest
$Z$ in the table with $22Z\ge8024$, i.e.\ $Z\ge364.7$. The candidates are 352 ($11\times32$, too small)
and 384 ($3\times128$), so $Z=384$ and $K=22\times384=8448$: the block is padded with
$8448-8024=424$ filler bits (known zeros, never transmitted, given infinite LLRs at the decoder).
The mother codeword has $68\times384=26\,112$ bits, of which the first $2Z=768$ are punctured,
leaving a circular buffer of $66\times384=25\,344$ bits. The first transmission carries
$E\approx8000/0.5=16\,000$ bits starting at RV0; the decoder sees a code of effective rate about 0.5,
built from roughly the first 23 rows of the base graph (the 7256 transmitted systematic bits plus
about $22.8Z$ parity bits).

Now a small block: $A=200$ bits at $R=0.3$. Since $A\le292$, \textbf{BG2}. The CRC is 16 bits
($B=216$). For $192<B\le560$ the table gives $K_b=8$, so the smallest $Z$ with $8Z\ge216$ is needed:
27 is not of the form $a\times2^j$, so $Z=28$ ($7\times4$). For BG2, $K=10Z=280$, with 64 filler bits;
the circular buffer has $50\times28=1400$ bits.
\end{worked}

\subsection{Other LDPC standards}

\begin{table}[htbp]\centering\small
\caption{LDPC codes in some important standards.}
\label{tab:ch15:ldpcstd}
\footnotesize
\begin{tabularx}{\linewidth}{@{}l l l >{\raggedright\arraybackslash}X@{}}
\toprule
Standard & Lengths $n$ & Rates & Structure and remarks\\
\midrule
5G NR data (TS 38.212) & $\le$26\,112 & 1/5--0.95 & two QC base graphs, 51 lifting sizes, rate compatible\\
Wi-Fi 802.11n/ac/ax & 648, 1296, 1944 & 1/2, 2/3, 3/4, 5/6 & QC, 24 block columns, $Z=27,54,81$\\
WiMAX 802.16e & 576--2304 & 1/2 to 5/6 & QC, $Z=24$ to 96\\
DVB-S2/S2X, DVB-T2, ATSC 3.0 & 16\,200, 64\,800 & 1/4 to 9/10 & IRA, 360-column groups, outer BCH\\
10GBASE-T (802.3an) & 2048 & 1723/2048 & RS-based construction, regular\\
CCSDS 131.0-B & 1280--32\,768 & 1/2--7/8 & AR4JA protographs; near-Earth (8160,7136)\\
ITU-T G.hn (G.9960) & 336--8640 & 1/2 to 5/6 & QC, home networking over wires\\
\bottomrule
\end{tabularx}
\end{table}

Table~\ref{tab:ch15:ldpcstd} collects the main deployments. Two deserve comment. The Wi-Fi code,
introduced as an option in 802.11n and now used for almost all high-throughput links, has three
lengths and four rates, each defined by a $24$-column base matrix; the rate-1/2 base matrix has 12
rows. It gains about 2\,dB over the $K=7$ convolutional code at the error rates Wi-Fi operates at,
which at high SNR is worth one or two MCS steps. The DVB-S2 code (ETSI EN 302 307-1), selected in
2003 from proposals by seven companies, is an irregular repeat--accumulate code of length 64\,800
with an outer BCH code correcting 8 to 12 errors. Its information part is built from groups of 360
columns that are cyclic shifts of each other, so a decoder can use 360 parallel processors; Chapter
\ref{ch:satellite} describes its place in the DVB-S2 physical layer. At the time it was the most
complex decoder ever standardised for a consumer product, and it operates roughly 0.6 to 1\,dB from
the constrained capacity at a packet error rate of $10^{-7}$.

\subsection{Error floors and trapping sets}

LDPC codes have error floors too, but their cause is different from the turbo-code case: not
low-weight codewords (good LDPC codes have large minimum distance) but \term{trapping sets}, small
subgraphs on which the iterative decoder gets stuck. An $(a,b)$ trapping set is a set of $a$ variable
nodes whose induced subgraph has $b$ odd-degree (unsatisfied) check nodes. If those $a$ bits are in
error, only $b$ checks complain, and if $b$ is small relative to $a$ the variables' other, satisfied
checks outvote the complaints; the decoder oscillates or converges to a wrong word that is not a
codeword. On the BEC the analogous objects are stopping sets. Trapping-set floors are sensitive to
the decoder's quantisation, saturation and scheduling, which makes them notoriously hard to predict
by simulation: an error rate of $10^{-12}$ at 10\,Gb/s is one error every 100 seconds, and verifying
it needs FPGA emulation of the decoder running for days.

Remedies are applied at every level. Code design avoids small trapping sets by maximising girth and
controlling the cycle structure (the approximate cycle extrinsic message degree, ACE, is one metric),
and by avoiding too many degree-2 variable nodes in cycles. Decoders use careful message
quantisation, post-processing that detects a stuck decoder and perturbs the LLRs of the bits
involved in unsatisfied checks, or \emph{backtracking}. Systems add an outer algebraic code with a
small correction capability, like DVB-S2's BCH code, which removes the few residual errors that a
trapping set leaves.

\subsection{LDPC decoders in hardware}

LDPC decoding is naturally parallel: every check node can be updated at once. Architectures span a
wide range:
\begin{itemize}
\item \textbf{Partially parallel, layered} decoders process one block row of a QC code per clock
cycle with $Z$ (or a multiple of $Z$) processing units. They dominate in 5G basebands and Wi-Fi.
\item \textbf{Fully parallel} decoders instantiate every node and wire the Tanner graph in silicon.
They reach the highest throughput but the routing congestion limits them to shortish codes;
10GBASE-T's 2048-bit code was an early example.
\item \textbf{Unrolled} decoders pipeline each iteration in its own hardware stage and have been
reported at throughputs of hundreds of Gb/s in research ASICs, at the cost of area.
\end{itemize}
Throughput scales as $T\approx(\text{parallelism}\times f_{\text{clk}}\times K)/(\text{cycles per
iteration}\times I)$, where $I$ is the number of iterations, so early termination (stop as soon as
the syndrome is zero) raises average throughput and cuts power. The 5G requirement of peak data rates
up to 20\,Gb/s on the downlink, combined with the need for LDPC's parallelism and early termination,
was a decisive argument in 2016.

\begin{pitfall}[Believing a waterfall that you have not simulated deep enough]
A code that looks excellent at a frame error rate of $10^{-3}$ can floor at $10^{-6}$. Monte Carlo
simulations that stop at 100 frame errors per point cannot see a floor below the last point, and
software decoders in double precision behave differently from 5-bit hardware. For any system with a
low target error rate---optical transport, storage, broadcast video---simulate the \emph{quantised}
decoder, with the exact schedule and the maximum iteration count, on hardware emulation, down to
at least a decade below the target, or use an outer code.
\end{pitfall}

%======================================================================
\section{Polar codes}
\label{sec:ch15:polar}
%======================================================================

Turbo and LDPC codes approach capacity in practice, but neither has a proof that it
\emph{achieves} capacity on general channels with an efficient decoder. (Irregular LDPC ensembles do
on the erasure channel, and spatially coupled ones do more generally, as Section~\ref{sec:ch15:sc}
explains.) In 2009 Erdal Ar\i kan published the first explicit construction of codes that provably
achieve the capacity of any binary-input memoryless symmetric channel with encoding and decoding
complexity $O(N\log N)$. The idea, \term{channel polarization}, is unlike anything else in coding
theory, and it is simple enough to explain completely with the erasure channel.

\begin{history}[From the cutoff rate to polarization]
Erdal Ar\i kan, a student of Gallager's at MIT in the 1980s and later a professor at Bilkent
University in Ankara, spent much of his career on sequential decoding and the cutoff rate $R_0$.
Pinsker and Massey had noticed that clever combinations of channels could push the sum of their
cutoff rates above the original; Ar\i kan was studying when, and how far, this ``boosting'' could
go. In his 2009 paper in the \emph{IEEE Transactions on Information Theory}, ``Channel polarization:
a method for constructing capacity-achieving codes for symmetric binary-input memoryless channels'',
he showed that a recursive combination of two channels into two new ones, applied $n$ times, drives
every synthetic channel towards either perfect or useless. The paper won the IEEE Information Theory
Society's best paper award, and its author the Shannon Award in 2018. In between, polar codes went
from a theoretical curiosity---their performance at practical lengths with the original decoder was
mediocre---to the control-channel code of 5G, a journey made possible by Tal and Vardy's list
decoder (2011) and an intensive industrial development effort led by Huawei.
\end{history}

\subsection{One step of polarization}

Take two independent uses of a binary-input channel $W$ and two input bits $u_1,u_2$. Instead of
sending them directly, send
\begin{equation}
x_1=u_1\oplus u_2,\qquad x_2=u_2,
\qquad\text{i.e.}\quad (x_1,x_2)=(u_1,u_2)\,\vect F,\quad
\vect F=\begin{bmatrix}1&0\\1&1\end{bmatrix}.
\label{eq:ch15:F}
\end{equation}
This invertible transform loses nothing: $I(U_1U_2;Y_1Y_2)=2I(W)$. Now decode in order. To decide
$u_1$, the receiver must work through $x_1=u_1\oplus u_2$ without knowing $u_2$: the effective channel
for $u_1$, written $W^-$, is \emph{worse} than $W$. Having decided $u_1$, the receiver sees $u_2$
twice---once directly through $y_2$ and once through $y_1$, since $x_1\oplus u_1=u_2$---so the
channel for $u_2$, written $W^+$, is \emph{better} than $W$. By the chain rule,
\begin{equation}
I(W^-)+I(W^+)=2I(W),\qquad I(W^-)\le I(W)\le I(W^+).
\label{eq:ch15:conserve}
\end{equation}
Capacity is conserved but redistributed.

For the BEC with erasure probability $\varepsilon$ the two new channels are again erasure channels.
$u_1$ is recovered only if both $y_1$ and $y_2$ arrive, so $W^-$ is a BEC with erasure probability
$1-(1-\varepsilon)^2=2\varepsilon-\varepsilon^2$; $u_2$ is lost only if both copies are erased, so
$W^+$ has erasure probability $\varepsilon^2$:
\begin{equation}
\varepsilon^-=2\varepsilon-\varepsilon^2,\qquad\varepsilon^+=\varepsilon^2.
\label{eq:ch15:becpol}
\end{equation}
For $\varepsilon=0.5$, one step produces channels with erasure probabilities 0.75 and 0.25.

\subsection{Recursion, polarization and the Kronecker construction}

Apply the step recursively: combine two copies of $W^-$ to get $W^{--}$ and $W^{-+}$, two copies of
$W^+$ to get $W^{+-}$ and $W^{++}$, and so on. After $n$ levels there are $N=2^n$ synthetic
\term{bit-channels} $W_N^{(i)}$, one for each sign pattern. On the BEC their erasure probabilities
follow from~\eqref{eq:ch15:becpol} exactly. Figure~\ref{fig:ch15_polarization}(a) shows the
result: level by level, the capacities spread towards 0 and 1. Ar\i kan proved that for any binary
symmetric memoryless channel, as $N\to\infty$, the fraction of bit-channels with capacity near 1
tends to $I(W)$ and the fraction with capacity near 0 to $1-I(W)$. The channels
\term{polarize}. The coding strategy follows: send information on the good bit-channels and fix the
bad ones to known values (zeros), the \term{frozen bits}. The rate approaches capacity.

\mdcfig{ch15_polarization}{Channel polarization on the BEC with erasure probability 0.5 (capacity
0.5). (a) Capacities $1-Z$ of the $2^n$ bit-channels at each level $n$ (horizontally jittered for
visibility). (b) Sorted capacities for $N=2^4$ to $2^{20}$: the curve approaches a step at 0.5, but
slowly. At $N=1024$ about a third of the channels are still ``mediocre'' (erasure probability between
$10^{-3}$ and $1-10^{-3}$); at $N=2^{20}$, about 5\%.}

The encoder is the $n$-fold Kronecker power of $\vect F$,
\begin{equation}
\vect x=\vect u\,\vect G_N,\qquad\vect G_N=\vect F^{\otimes n},
\label{eq:ch15:polarenc}
\end{equation}
where $\vect u$ has the information bits in the positions of the good bit-channels (the
\term{information set} $\mathcal A$) and zeros elsewhere. (Ar\i kan's original definition includes a
bit-reversal permutation, $\vect G_N=\vect B_N\vect F^{\otimes n}$; 5G NR and \texttt{commlib} use
$\vect F^{\otimes n}$ alone, which only relabels the bit-channels.) The product can be computed with
$\frac N2\log_2N$ XORs in the butterfly network of Figure~\ref{fig:ch15_polar_encoder}, the same
structure as the FFT.

\begin{figure}[htbp]\centering
\begin{tikzpicture}[x=1.45cm,y=0.52cm,
  xo/.style={circle,draw=navy,thick,inner sep=0pt,minimum size=3.6mm,fill=white},
  dt/.style={circle,fill=navy,inner sep=0pt,minimum size=1.6mm}]
\foreach \i in {0,...,7}{
  \draw[navy!50,thick] (0,-\i) -- (4,-\i);
}
\foreach \i/\fz in {0/1,1/1,2/1,3/0,4/1,5/0,6/0,7/0}{
  \ifnum\fz=1
    \node[anchor=east,font=\sffamily\scriptsize,gray] at (-0.05,-\i) {$u_{\i}=0$ (frozen)};
  \else
    \node[anchor=east,font=\sffamily\scriptsize,navy] at (-0.05,-\i) {$\mathbf{u_{\i}}$ (information)};
  \fi
  \node[anchor=west,font=\sffamily\scriptsize] at (4.05,-\i) {$x_{\i}\to W$};
}
% stage spans 4, 2, 1 at x = 1, 2, 3
\foreach \s/\x in {4/1,2/2,1/3}{
  \foreach \i in {0,...,7}{
    \pgfmathtruncatemacro{\m}{mod(\i,2*\s)}
    \ifnum\m<\s
      \pgfmathtruncatemacro{\j}{\i+\s}
      \draw[navy,thick] (\x,-\i) -- (\x,-\j);
      \node[xo] at (\x,-\i) {\tiny$+$};
      \node[dt] at (\x,-\j) {};
    \fi
  }
}
\node[font=\sffamily\scriptsize,navy] at (1,0.9) {span 4};
\node[font=\sffamily\scriptsize,navy] at (2,0.9) {span 2};
\node[font=\sffamily\scriptsize,navy] at (3,0.9) {span 1};
\end{tikzpicture}
\caption{The polar encoder $\vect x=\vect u\vect F^{\otimes3}$ for $N=8$ as a butterfly network of
$\frac N2\log_2N=12$ XORs. Each stage applies the $2\times2$ kernel $\vect F$ to pairs of lines a
fixed span apart. With $K=4$ and a design for the BEC(0.5), the information set is
$\mathcal A=\{3,5,6,7\}$, the four bit-channels with the smallest erasure probabilities
(their erasure probabilities are 0.32, 0.19, 0.12 and 0.004; the best frozen channel, $u_4$, has 0.68).}
\label{fig:ch15_polar_encoder}
\end{figure}

\subsection{Choosing the frozen set}

Constructing a polar code means ranking the $N$ bit-channels by reliability for the channel at hand
and freezing the worst $N-K$. On the BEC the recursion~\eqref{eq:ch15:becpol} gives the erasure
probabilities exactly. For other channels the bit-channel output alphabets grow exponentially and
exact computation is impossible, so several approximations are used:
\begin{itemize}
\item \textbf{Bhattacharyya bounds.} Track $Z(W)=\sum_y\sqrt{W(y|0)W(y|1)}$, which upper-bounds the
error probability, with $Z^-\le2Z-Z^2$ and $Z^+=Z^2$ (exact on the BEC). Simple, and the method in
\texttt{commlib.PolarCode}, but pessimistic for the AWGN channel.
\item \textbf{Density evolution with quantisation} (Tal and Vardy, 2013) computes upper and lower
bounds on each bit-channel's error probability by merging output symbols; it is accurate but slow.
\item \textbf{Gaussian approximation} (Trifonov, 2012) tracks the mean LLR of each bit-channel as
in~\eqref{eq:ch15:ga}: $\mu^+=2\mu$ for the good channel and
$\mu^-=\phi_G^{-1}\big(1-(1-\phi_G(\mu))^2\big)$ for the bad. It is fast and accurate, and used for the
simulations in this chapter.
\item \textbf{A fixed reliability sequence.} All of the above depend on a design SNR. 5G NR instead
defines a single ordering $\mathcal Q_0^{1023}$ of the 1024 bit-channel indices from least to most
reliable (TS 38.212, Table 5.3.1.2-1). It is \emph{nested}: the ordering for any $N\le1024$ is the
subsequence of indices smaller than $N$. Every encoder and decoder stores just this one table. The
sequence came from extensive optimisation; it is close to the ``polarization weight'' ordering, in
which bit-channel $i$ with binary digits $b_j$ gets weight $\sum_jb_j2^{j/4}$.
\end{itemize}

\subsection{Successive-cancellation decoding}

The decoder implied by the construction decodes $u_0,u_1,\dots,u_{N-1}$ in order, each decision
using the channel outputs and all previous decisions: \term{successive cancellation} (SC). Frozen
bits are set to zero; an information bit is set by the sign of its LLR. The LLRs are computed by
running the butterfly backwards with two operations that are just the soft versions of the two
channels in~\eqref{eq:ch15:F}. For a pair of LLRs $(a,b)$ at the outputs of a kernel,
\begin{align}
f(a,b)&=a\boxplus b\approx\operatorname{sign}(a)\operatorname{sign}(b)\min(|a|,|b|)
&&\text{(LLR of the upper input, through the XOR)},\nonumber\\
g(a,b,\hat u)&=b+(1-2\hat u)\,a&&\text{(LLR of the lower input, given the upper)}.
\label{eq:ch15:fg}
\end{align}
The upper input (like $u_1$ in~\eqref{eq:ch15:F}) is seen through an XOR, hence the box-plus;
the lower input, once the upper is known, is seen twice, hence the sum, with the sign of $a$
flipped if the upper bit was 1. SC decoding is a depth-first traversal of a binary tree
(Figure~\ref{fig:ch15_sc_tree}): each node receives a vector of LLRs from its parent, computes $f$ for
its left child, waits for the left child's hard decisions (partial sums) $\beta_L$, computes $g$ for its
right child, and returns $(\beta_L\oplus\beta_R,\beta_R)$ to its parent. The complexity is
$O(N\log N)$, exactly like the encoder.

\begin{figure}[htbp]\centering
\begin{tikzpicture}[level distance=1.05cm,
  level 1/.style={sibling distance=4.6cm},
  level 2/.style={sibling distance=2.3cm},
  level 3/.style={sibling distance=1.15cm},
  every node/.style={circle,draw=navy,thick,minimum size=5mm,inner sep=0pt,font=\sffamily\tiny,fill=navy!8},
  edge from parent/.style={draw=navy!60,thick}]
\node {$N{=}8$}
  child {node[fill=gray!20] {R0}
    child {node[fill=gray!20] {} child {node[fill=white] {$u_0$}} child {node[fill=white] {$u_1$}}}
    child {node {} child {node[fill=white] {$u_2$}} child {node[fill=navy!50,text=white] {$u_3$}}}
    edge from parent node[draw=none,fill=none,left,font=\sffamily\scriptsize,accent] {$f$}}
  child {node {}
    child {node {} child {node[fill=white] {$u_4$}} child {node[fill=navy!50,text=white] {$u_5$}}}
    child {node[fill=navy!50,text=white] {R1} child {node[fill=navy!50,text=white] {$u_6$}} child {node[fill=navy!50,text=white] {$u_7$}}}
    edge from parent node[draw=none,fill=none,right,font=\sffamily\scriptsize,practc] {$g$}};
\node[draw=none,fill=none,font=\sffamily\scriptsize,anchor=west,align=left] at (4.3,-0.2)
  {$\alpha$ (LLRs) flow down:\\ left child gets $f(\alpha_{\text{top}},\alpha_{\text{bot}})$,\\
   right child $g(\alpha_{\text{top}},\alpha_{\text{bot}},\beta_L)$;\\
   $\beta$ (partial sums) flow up:\\ $(\beta_L\oplus\beta_R,\;\beta_R)$};
\end{tikzpicture}
\caption{Successive-cancellation decoding as a depth-first traversal of a binary tree, for the
$N=8$ code of Figure~\ref{fig:ch15_polar_encoder}. Leaves are bit-channels: white frozen, dark
information. A subtree whose leaves are all frozen (a rate-0 node, R0) needs no computation; one
whose leaves are all information (rate-1, R1) can be decided at once by hard decisions on its input
LLRs. Fast-SSC decoders exploit these and other special nodes to shorten the traversal greatly.}
\label{fig:ch15_sc_tree}
\end{figure}

\begin{worked}[Decoding a length-4 polar code by hand]
For the BEC(0.5) design with $N=4$, the bit-channel erasure probabilities are $0.94$, $0.56$, $0.44$ and
$0.06$, so a $K=2$ code freezes $u_0=u_1=0$ and puts information on $u_2,u_3$. Send $u_2=u_3=1$; with
$\vect G_4=\vect F^{\otimes2}$ the codeword is $\vect x=(0,1,0,1)$, transmitted as $(+1,-1,+1,-1)$.
The received LLRs are $\vect L=(+1.5,-0.8,-0.4,-2.0)$: note that the third bit arrived with the wrong
sign. Decode with min-sum SC.

The root splits $\vect L$ into halves $(1.5,-0.8)$ and $(-0.4,-2.0)$. Left child:
$f(1.5,-0.4)=-0.4$ and $f(-0.8,-2.0)=+0.8$. Inside it, $u_0$ gets $f(-0.4,0.8)=-0.4$ but is frozen,
so $\hat u_0=0$; $u_1$ gets $g(-0.4,0.8,0)=0.4$, frozen, $\hat u_1=0$. The left child returns partial
sums $\beta_L=(0,0)$. Right child: $g(1.5,-0.4,0)=1.1$ and $g(-0.8,-2.0,0)=-2.8$. Inside it, $u_2$ gets
$f(1.1,-2.8)=-1.1$, so $\hat u_2=1$; $u_3$ gets $g(1.1,-2.8,1)=-2.8-1.1=-3.9$, so $\hat u_3=1$. Both
information bits are correct despite the channel error: the frozen bits and the structure of
$\vect G_4$ corrected it.
\end{worked}

SC decoding achieves capacity as $N\to\infty$, with block error probability falling roughly as
$2^{-\sqrt N}$. But it is weak at practical lengths (Figure~\ref{fig:ch15_polar_decoders}): the
mediocre bit-channels that have not yet polarized make errors, and an error in an early bit is never
corrected---it propagates through every later $g$ computation. At $N=1024$, SC needs about 1.2\,dB
more than the best decoder below.

\subsection{List decoding and the CRC}

Ido Tal and Alexander Vardy's \term{successive-cancellation list} (SCL) decoder (2011) keeps up to
$L$ candidate paths. At each information bit every path is extended both ways, each extension is
scored by a \emph{path metric}, and the $L$ best survive. With LLRs the path metric is updated, to an
excellent approximation, by
\begin{equation}
\mathrm{PM}\leftarrow\mathrm{PM}+\begin{cases}|\alpha_i|&\text{if the decision }\hat u_i\text{
disagrees with the sign of its LLR }\alpha_i,\\0&\text{otherwise,}\end{cases}
\label{eq:ch15:pm}
\end{equation}
so a path pays a penalty whenever it goes against the evidence (frozen bits included: a frozen zero
with a strongly negative LLR penalises the path). At the end, the path with the smallest metric wins.
SCL with $L=32$ comes close to maximum-likelihood decoding of the polar code---and that exposed the
next problem: the polar code itself has a small minimum distance, so even ML decoding is mediocre
(the ``SCL, $L=8$'' curve of Figure~\ref{fig:ch15_polar_decoders} floors for exactly this reason).

The fix is to append a CRC to the information bits before polar encoding, and to choose, at the end,
the most likely path \emph{that passes the CRC}: \term{CRC-aided SCL} (CA-SCL; Niu and Chen, 2012).
The CRC acts both as an outer code that increases the minimum distance and as a genie that picks the
right path from the list. Polar codes with CA-SCL decoding are the strongest known practical codes at
short block lengths. Figure~\ref{fig:ch15_polar_decoders} shows the progression at $N=1024$, $K=512$:
CA-SCL with $L=32$ reaches a frame error rate of $10^{-2}$ about 1.2\,dB before SC and within about
0.6\,dB of the normal-approximation bound.

\mdcfig{ch15_polar_decoders}{Frame error rate of rate-1/2 polar codes with $N=1024$ (GA construction
at 2\,dB, min-sum $f$ function) on the BPSK/AWGN channel: SC, SC-list with $L=8$ and no CRC, and
CRC-aided SCL with the 5G CRC-11 (so $K=512$ message bits plus 11 CRC bits on 523 bit-channels) and
$L=8$ or 32. Dashed: the normal-approximation bound for any $(1024,512)$ code. At least 60 frame
errors or 1500--4000 frames per point.}

\begin{pitfall}[The CRC is not free]
A CRC of $r$ bits that is used to choose among $L$ list candidates both improves performance and
weakens error \emph{detection}. Each wrong candidate that is tested has roughly a $2^{-r}$ chance of
passing, so the probability of an undetected error grows to about $L\cdot2^{-r}$. With $L=8$ and an
11-bit CRC that is about $4\times10^{-3}$ of the decoding failures---far too high for a control
channel that must almost never deliver a false command. 5G handles this by using a 24-bit CRC for
downlink control (giving $L\cdot2^{-24}$, negligible) and by counting the bits the list ``spends''
when it sets the CRC length. A decoder designer who increases $L$ without accounting for this trades
undetected errors for frame errors.
\end{pitfall}

\subsection{Polar codes in 5G NR}

5G NR uses polar codes for the broadcast channel and for control information in both directions
(TS 38.212, Sections 5.3.1, 7.1 and 7.3). Table~\ref{tab:ch15:nrpolar} summarises the uses.

\begin{table}[htbp]\centering\small
\caption{Polar codes in 5G NR (3GPP TS 38.212). $A$ is the payload, $K$ the number of bits entering
the polar encoder, $N\le2^{n_{\max}}$ the mother code length.}
\label{tab:ch15:nrpolar}
\footnotesize
\begin{tabularx}{\linewidth}{@{}l l l l >{\raggedright\arraybackslash}X@{}}
\toprule
Channel & Payload & CRC & $n_{\max}$ & Remarks\\
\midrule
PBCH (broadcast) & $A=32$ & 24 (CRC-24C) & 9 & $K=56$, $E=864$, $N=512$\\
PDCCH (DCI) & up to $\approx$140 & 24 (CRC-24C) & 9 & distributed CRC; last 16 CRC bits masked with the RNTI\\
UCI, $12\le A\le19$ & 12--19 & 6 + 3 parity-check bits & 10 & PC-polar; on PUCCH or PUSCH\\
UCI, $A\ge20$ & 20--1706 & 11 & 10 & segmented into two blocks for large $A$ and $E$\\
\bottomrule
\end{tabularx}
\end{table}

Several details deserve comment.
\begin{itemize}
\item \textbf{Very small payloads} ($A\le11$ bits of uplink control) do not use polar codes at all:
1 and 2 bits use repetition and simplex codes, 3 to 11 bits a $(32,A)$ Reed--Muller code.
\item \textbf{Distributed CRC.} For downlink control the CRC bits are interleaved among the information
bits so that each CRC bit depends only on earlier information bits. An SCL decoder can then check
parts of the CRC as it goes and abandon a candidate PDCCH early. That matters because a terminal
must attempt blind decoding of dozens of PDCCH candidates in every slot (up to 44 per slot at
15\,kHz subcarrier spacing), most of which contain nothing for it.
\item \textbf{Parity-check bits.} For short uplink payloads, three parity-check bits computed by a
small shift register over the information bits are placed on specific frozen positions, giving the
list decoder something to check without the cost of a longer CRC.
\item \textbf{Rate matching.} The coded bits pass through a 32-segment sub-block interleaver into a
circular buffer. If the number of transmitted bits $E$ is less than $N$, bits are either
\emph{punctured} (removed from the start; used at low rates, $K/E\le7/16$) or \emph{shortened}
(removed from the end, with the corresponding bit-channels frozen so that the receiver knows them;
used at higher rates). If $E>N$, bits are repeated. The frozen set is adjusted to account for the
punctured or shortened positions. $N$ is chosen as the smallest power of two that keeps the mother
code rate at or below about $1/8$ relative to $E$, capped by $n_{\max}$, with a rule that prefers the
smaller $N$ when only a little puncturing is needed.
\item \textbf{Interleaving.} Uplink control uses a triangular channel interleaver so that the coded
bits of a high-order constellation symbol are spread across the codeword.
\end{itemize}

\begin{inpractice}[Polar decoders in silicon]
SC decoding is sequential in principle, but the tree structure allows a great deal of parallelism.
Fast simplified SC decoders (Sarkis, Giard, Vardy, Thibeault and Gross, 2014) recognise rate-0,
rate-1, repetition and single-parity-check subtrees and decode each in one or two clock cycles,
reducing the latency by an order of magnitude. List decoders add the cost of sorting $2L$ path metrics
and copying partial sums; $L=8$ is the usual choice in 5G terminals. Because 5G control blocks are
short (at most 512 bits downlink and 1024 uplink), a CA-SCL decoder is a small block compared with the
LDPC decoder that shares the modem with it. Fully unrolled and pipelined SC decoders for fixed codes
have been reported with throughputs of hundreds of Gb/s in research ASICs, which has made polar codes
candidates for very high-rate links too.
\end{inpractice}

%======================================================================
\section{Comparing the families}
\label{sec:ch15:compare}
%======================================================================

\subsection{Short blocks}

Chapter~\ref{ch:infotheory} derived the \term{normal approximation} for the largest number of
information bits $k$ that any code of length $n$ can carry at block error probability $\epsilon$,
\begin{equation}
k\approx nC-\sqrt{nV}\,\Q^{-1}(\epsilon)+\tfrac12\log_2n,
\label{eq:ch15:na}
\end{equation}
where $C$ is the capacity and $V$ the channel dispersion (both for the BPSK-input AWGN channel
here, computed by numerical integration). Inverting it gives the smallest SNR at which a code of
given $n$, $k$ and $\epsilon$ can exist: the yardstick for real codes. For $(256,128)$ at
$\epsilon=10^{-3}$ it is about 1.9\,dB, 1.7\,dB above the asymptotic limit.

\mdcfig{ch15_short_compare}{Four code families at $k=128$ information bits and rate about 1/2:
polar (256,128) with CA-SCL ($L=8$ and 32, CRC-11 included in the 128 bits' overhead), a (3,6)-regular
PEG LDPC code of length 256 decoded by sum-product (100 iterations), an LTE-type turbo code with
$K=128$ punctured to rate $\approx$0.48 (QPP $f_1=15$, $f_2=32$, 8 log-MAP iterations), and the $K=7$
convolutional code with Viterbi decoding, all with $n\approx256$--268 coded bits. Dashed: the
normal-approximation bound. At least 60 frame errors or up to 8000 frames per point.}

Figure~\ref{fig:ch15_short_compare} compares the families at $k=128$, a typical size for control
information. At a frame error rate of $10^{-3}$, polar CA-SCL with $L=32$ needs about 2.6\,dB and with
$L=8$ about 2.9\,dB; the turbo code about 3.0\,dB; the regular LDPC code about 3.3\,dB; and the
convolutional code about 4.0\,dB. The best is within about 0.7\,dB of the bound. A well-designed
irregular LDPC code such as NR's BG2 does better than the regular code shown, but at these lengths
it does not catch CA-SCL polar codes, and LDPC codes at short lengths suffer from the cycles that a
short, sparse graph cannot avoid. Tail-biting convolutional codes with ML (wrap-around Viterbi)
decoding---LTE's control-channel code---are surprisingly competitive at the very shortest lengths, a
few tens of bits, but fall behind as $k$ grows.

\mdcfig{ch15_gap_length}{Required $\ebno$ for a frame error rate of $10^{-2}$ as a function of the
number of information bits, for BPSK on the AWGN channel. Solid lines: normal approximation for rates
1/3, 1/2 and 3/4; dotted lines: the corresponding asymptotic (Shannon) limits. Markers: codes
simulated for this chapter---the rate-1/3 LTE-type turbo codes of Figure~\ref{fig:ch15_turbo_length},
the four rate-1/2 codes of Figure~\ref{fig:ch15_short_compare}, and the $(3,6)$ QC-LDPC code of
Figure~\ref{fig:ch15_ldpc_decoders}.}

\subsection{Long blocks}

Figure~\ref{fig:ch15_gap_length} puts the comparison on one chart. The finite-length penalty is
severe at short lengths---about 2\,dB at $k=100$ for rate 1/2---and falls roughly as $1/\sqrt k$; at
$k=10\,000$ it is a few tenths of a decibel. Good practical codes track the bound at a roughly
constant distance of half a decibel to a decibel. At long block lengths, irregular LDPC and turbo
codes are both within about a decibel of capacity at practical error rates, and the choice between
them is made on other grounds: decoder throughput and parallelism (LDPC wins), error floors (LDPC
with good design and an outer code wins), flexibility of length and rate (both can be made flexible;
NR LDPC shows how), latency and energy per bit (LDPC with layered decoding and early termination),
and intellectual property and implementation maturity, which in 2016 were real considerations.

\begin{history}[The 5G coding decision, 2016]
3GPP's study of the New Radio began in 2016 with three candidates for the enhanced mobile broadband
(eMBB) data channels: an evolution of the LTE turbo code, LDPC codes and polar codes. The debate in
the RAN1 working group was technical and commercial at once: turbo codes had a decade of deployment;
LDPC codes had the throughput and parallelism needed for multi-gigabit peak rates and were backed by
Qualcomm, Samsung, Nokia, Ericsson, Intel and others; polar codes were championed above all by
Huawei, which had invested heavily in them. At RAN1 meeting \#86bis in Lisbon in October 2016,
LDPC was agreed for the eMBB data channels. At meeting \#87 in Reno, Nevada, in November 2016, after a
long session, polar codes were agreed for eMBB control channels in both directions (with the
broadcast channel following), and the turbo code was dropped. The split matched the technical
picture---LDPC for long blocks and high throughput, polar CA-SCL for short blocks---and also a
compromise among the companies. It was the first time a code family invented in the twenty-first
century had been chosen for a global cellular standard, eight years after its invention.
\end{history}

\begin{table}[htbp]\centering\small
\caption{The three families compared, as used in practice.}
\label{tab:ch15:compare}
\footnotesize
\begin{tabularx}{\linewidth}{@{}l >{\raggedright\arraybackslash}X >{\raggedright\arraybackslash}X >{\raggedright\arraybackslash}X@{}}
\toprule
 & Turbo & LDPC & Polar\\
\midrule
Construction & two RSCs + interleaver & sparse $\vect H$, protograph + lifting & Kronecker kernel + frozen set\\
Decoder & iterative BCJR (log-MAP) & BP / layered min-sum & SC, CA-SCL\\
Parallelism & windows, limited & very high & tree-level; good for short blocks\\
Best regime & medium blocks, low rates & long blocks, high rates and throughput & short blocks\\
Error floor & interleaver-dependent & trapping sets; outer code if needed & low with CRC; ML-like\\
Rate compatibility & puncturing & extension parity (NR) & puncturing/shortening\\
Key standards & 3G, 4G, CCSDS, DVB-RCS & 5G data, Wi-Fi, DVB-S2, 10GBASE-T & 5G control and PBCH\\
\bottomrule
\end{tabularx}
\end{table}

%======================================================================
\section{Hybrid ARQ and rate-compatible codes}
\label{sec:ch15:harq}
%======================================================================

A wireless link adapts its modulation and coding scheme to a channel estimate that is always a little
out of date (Chapter~\ref{ch:channels}). Some blocks will fail. \term{Hybrid automatic repeat
request} (HARQ) combines forward error correction with retransmission: if the CRC fails, the receiver
sends a negative acknowledgement, the transmitter sends more, and the receiver \emph{combines} the
new transmission with the soft information it kept from the old one. There are two ways to combine.
\begin{itemize}
\item \textbf{Chase combining} (CC): the retransmission is identical, and the receiver adds the LLRs
of the two copies. This is maximal-ratio combining: the SNR adds, but the code rate is unchanged.
\item \textbf{Incremental redundancy} (IR): the retransmission carries \emph{new} parity bits, so the
receiver ends up with a longer, lower-rate codeword. Information-theoretically, the mutual
information of the transmissions adds, which is always at least as good as adding SNRs, and much
better at high SNR, where the logarithm of an SNR sum grows slowly.
\end{itemize}
IR needs a \term{rate-compatible} code family: a mother code of low rate whose higher-rate members
are obtained by puncturing, such that the bits sent in the first transmission are a subset of the
mother codeword and later transmissions reveal more of it. Turbo codes provide it by puncturing the
parity streams; the NR LDPC base graphs provide it by construction (each extension row adds a parity
bit); polar codes are harder, and NR does not use HARQ for polar-coded control channels.

The circular buffer is the mechanism. The coded bits are written once into a buffer, and each
transmission reads a contiguous segment starting at its redundancy version's position, wrapping
round as necessary. In NR the usual redundancy version sequence is 0, 2, 3, 1: RV0 contains mostly
systematic bits (which gives the best first-transmission performance), RV2 and RV3 carry fresh parity,
and RV1 is sent last. LTE and NR terminals run several HARQ processes in parallel (8 in LTE FDD, up to 16
in NR) so that the link stays busy while each process waits for its acknowledgement.

\mdcfig{ch15_harq}{Throughput of HARQ on a Rayleigh block-fading channel (independent fade per
transmission), with a first-transmission rate of 4\,bits per channel use and up to four transmissions,
using an idealised capacity-based decoding model: a block is decoded once the accumulated mutual
information (IR) or the information at the accumulated SNR (CC) exceeds the rate. Plain ARQ discards
failed transmissions. 40\,000 channel realisations per point.}

Figure~\ref{fig:ch15_harq} shows the difference on a fading channel. At an average SNR of 10\,dB,
plain ARQ delivers about 0.9 bits per channel use, Chase combining about 1.6 and incremental redundancy
about 2.0. In effect, IR lets the transmitter pick an aggressive initial rate and lets the channel
decide, after the fact, how much redundancy was needed: it achieves a rate close to the mutual
information of the realised channel without knowing it in advance.

\begin{worked}[Chase combining versus incremental redundancy]
A codeword carrying 1.5 bits per transmitted symbol fails on its first transmission at an SNR of
2\,dB. The retransmission arrives at the same SNR. Using the idealised rule that a codeword is decoded
when the available mutual information exceeds the information it carries, does either scheme
succeed?

At 2\,dB ($\gamma=1.58$), one transmission provides $\log_2(1+1.58)=1.37$ bits per symbol, less than
1.5: failure. With \textbf{Chase combining} the two identical copies add their SNRs to $3.17$, giving
$\log_2(1+3.17)=2.06$ bits per symbol position, more than 1.5: success. With \textbf{incremental
redundancy} each symbol position of the original codeword now has two transmitted symbols carrying
different parity, so the available information is $2\times1.37=2.74$ bits per original position,
against the 1.5 required: success with a much larger margin, which a real system could have used by
starting at a higher rate. At low SNR the two schemes are nearly equivalent, because
$\log_2(1+2\gamma)\approx2\log_2(1+\gamma)$ when $\gamma\ll1$; at high SNR incremental redundancy is
far better, since doubling the SNR adds only one bit, whereas doubling the length doubles the bits.
\end{worked}

%======================================================================
\section{Spatially coupled LDPC codes}
\label{sec:ch15:sc}
%======================================================================

For a regular LDPC ensemble, BP decoding has a threshold strictly below what the same code could
achieve with optimal (MAP) decoding: for the $(3,6)$ ensemble on the BEC, 0.4294 against 0.4881.
Irregular designs close the gap to capacity by shaping the degree distribution, but they are
delicate: degree-2 nodes cause error floors, high degrees cause long iterations. In 1999 Alberto
Jim\'enez Felstr\"om and Kamil Zigangirov introduced \emph{LDPC convolutional codes}, and in 2011
Shrinivas Kudekar, Tom Richardson and R\"udiger Urbanke explained why they are so good, in one of the
most striking results of modern coding theory.

A \term{spatially coupled} LDPC code is built from $L$ copies of a protograph placed side by side in a
chain, with the edges of each copy spread over its neighbours within a window of width $w$. The
codes at the two ends of the chain have more checks than variables (lower rate), and that small,
local excess of information acts as a seed. In density evolution, the bits at the ends are decoded
first, and a \emph{decoding wave} then travels inward along the chain
(Figure~\ref{fig:ch15_spatial_coupling}a), decoding the middle even at erasure probabilities where an
uncoupled code would be stuck. Kudekar, Richardson and Urbanke proved \term{threshold saturation}: as
$L$ and $w$ grow, the BP threshold of the coupled ensemble converges to the MAP threshold of the
underlying uncoupled ensemble. Since the MAP thresholds of regular ensembles approach capacity as the
degrees grow, coupled regular codes achieve capacity on the BEC, and (as later proved) on all binary
memoryless symmetric channels, with BP decoding.

\mdcfig{ch15_spatial_coupling}{Spatial coupling on the BEC, by density evolution of the $(3,6,L,w=3)$
coupled ensemble. (a) Erasure probability along a chain of $L=64$ positions at $\varepsilon=0.47$
(above the uncoupled BP threshold), plotted every 80 iterations from dark to light: the ends decode
first and two waves move inward. (b) The coupled BP threshold versus chain length, compared with the
uncoupled BP threshold, the MAP threshold of the $(3,6)$ ensemble and the Shannon limit $1-R_L$ of the
coupled code, whose rate $R_L$ suffers a loss of order $1/L$ from the termination.}

The price is a rate loss of order $w/L$ from the termination, which vanishes as the chain grows
(Figure~\ref{fig:ch15_spatial_coupling}b: already at $L=16$ the threshold has saturated at 0.488), and
a long effective block length. Coupled codes are decoded with a \emph{sliding window} that processes a
few positions at a time, so the latency is set by the window, not the chain. They are attractive where
streams are long and continuous: optical transport, where several high-performance proposals are
spatially coupled or ``convolutional'' in structure, and the staircase codes of the next section,
which are a hard-decision cousin of the same idea.

\begin{keyidea}[Threshold saturation]
A regular LDPC code is held back not by its structure but by the inability of BP to get started.
Give it a small amount of side information at a boundary and couple its copies so that knowledge can
propagate, and BP achieves what an optimal decoder could. The wave picture---a front of known bits
travelling into a sea of unknown ones---is worth remembering; it reappears in coded random access and
in compressed sensing.
\end{keyidea}

%======================================================================
\section{Beyond radio: optical fibre and storage}
\label{sec:ch15:optics}
%======================================================================

\subsection{Optical transport}

Optical links have requirements unlike radio's: output bit error rates of $10^{-15}$ or below,
line rates of 100\,Gb/s to more than 1\,Tb/s per wavelength, and decoder power budgets of a few watts.
The history of optical FEC is a sequence of generations (Chapter~\ref{ch:wireline} gives the system
context):
\begin{itemize}
\item \textbf{First generation:} the Reed--Solomon RS(255,239) code of ITU-T G.709, with 6.7\%
overhead and hard-decision decoding, giving about 6\,dB of net coding gain at $10^{-15}$.
\item \textbf{Second generation:} concatenated and product-like hard-decision codes (ITU-T G.975.1),
gaining 8 to 9\,dB at the same overhead.
\item \textbf{Staircase codes} (Smith, Kschischang and co-workers, 2012) arrange the data in a chain
of square blocks, each block's rows forming BCH codewords together with the transposed previous
block. Each bit is protected by two component codes, as in a product code, and decoding proceeds
iteratively along the chain, a hard-decision cousin of spatial coupling. With 6.7\% overhead
(rate 239/255) they gain about 9.4\,dB at $10^{-15}$, within roughly 0.5\,dB of the hard-decision
capacity, and are standardised in ITU-T G.709.2.
\item \textbf{Soft-decision FEC} for coherent transceivers (Chapter~\ref{ch:wireline}): the OIF 400ZR
standard uses \emph{CFEC}, a concatenation of an inner soft-decision Hamming code and an outer
staircase code (about 15\% overhead), and the open \emph{oFEC} of OpenROADM and OpenZR+ uses a
block-based, soft-decision decoded design that tolerates a pre-FEC BER of roughly $2\times10^{-2}$.
Long-haul proprietary transceivers use soft-decision LDPC or spatially coupled codes with 15 to 27\%
overhead.
\end{itemize}
The engineering constraint is energy: at 400\,Gb/s, a decoder that spends 10\,pJ per bit dissipates
4\,W, a large fraction of a pluggable module's budget. That pushes optical FEC towards hard or
``soft-aided'' decoding of algebraic component codes, where it can be had.

\subsection{Flash memory and hard drives}

NAND flash stores data as charge levels in cells: 2 levels per cell for SLC, 4 for MLC, 8 for TLC and
16 for QLC. As the number of levels grew and the cells shrank, raw bit error rates rose to around
$10^{-3}$ or worse near end of life, and the BCH codes of early solid-state drives ran out of strength.
From the early 2010s, SSD controllers adopted LDPC codes, typically quasi-cyclic with codeword lengths
of a few kilobytes and rates around 0.9. Soft information is obtained by reading the page several
times with shifted reference voltages: each extra read refines the LLR of each cell's bits, at a cost
in latency, so controllers use hard-decision (bit-flipping or min-sum) decoding first and fall back
to soft reads only when it fails. Hard-disk drives made the same transition earlier: after PRML and
NPML detection (Chapter~\ref{ch:baseband}), read channels moved in the late 2000s to iterative
detection, exchanging soft information between a soft-output detector for the recording channel and
an LDPC decoder---turbo equalization (Chapter~\ref{ch:equalization}) inside every hard drive.

\begin{inpractice}[Where the iterative decoders are]
A modern smartphone contains, at least, a 5G LDPC decoder and polar decoder, an LTE turbo decoder, a
Wi-Fi LDPC decoder and the LDPC decoder in its flash controller. A data-centre server adds LDPC in its
SSDs and soft-decision FEC in its optical transceivers. Satellite TV receivers decode DVB-S2 LDPC
codes; deep-space ground stations decode CCSDS turbo and LDPC codes. The ideas of this chapter run
on billions of chips, and most of them use the min-sum check node.
\end{inpractice}

%======================================================================
\section{Outlook}
\label{sec:ch15:outlook}
%======================================================================

Capacity-approaching codes are, for long blocks, a solved problem; the open questions are elsewhere.
\begin{itemize}
\item \textbf{Short blocks and ultra-reliability.} Machine-type and ultra-reliable low-latency
communication (URLLC) need blocks of tens to hundreds of bits at error rates of $10^{-5}$ to
$10^{-9}$, where the finite-length penalty dominates. CA-SCL polar codes with large lists,
polarization-adjusted convolutional (PAC) codes (Ar\i kan, 2019), and short codes with near-ML
decoding are all close to the normal-approximation bound there.
\item \textbf{Universal decoders.} \term{Guessing random additive noise decoding} (GRAND; Duffy, Li
and M\'edard, 2019) inverts the usual approach: it guesses the noise pattern, from most likely to
least likely, subtracts it, and checks whether the result is a codeword. It works for \emph{any}
short, high-rate code (including CRCs alone), is ML for well-ordered guesses, and has spawned
soft-input variants (ORBGRAND) and hardware prototypes.
\item \textbf{Learned decoders.} Neural BP (Nachmani, Be'ery and Burshtein, 2016) unrolls belief
propagation into a network and learns a weight for every edge and iteration, recovering some of the
loss caused by short cycles in short codes---effectively a generalisation of normalised min-sum.
Fully learned encoders and decoders have reproduced classical performance in some settings, and
learned components are being studied for 6G (see Chapter~\ref{ch:sixg}).
\item \textbf{Coding with modulation and shaping.} Probabilistic amplitude shaping
(Chapter~\ref{ch:modulation}) combined with LDPC codes approaches the full Gaussian capacity rather than
the uniform-QAM capacity, and is already in optical transceivers and the newest Wi-Fi and DVB
proposals.
\end{itemize}

%======================================================================
\section{Summary}
%======================================================================

\begin{itemize}
\item Modern decoders trade in log-likelihood ratios. Independent evidence adds; the XOR of bits
combines by the box-plus (tanh) rule, well approximated by sign-times-minimum. Only
\emph{extrinsic} information may be passed between cooperating decoders.
\item The BCJR algorithm computes exact bit APPs on a trellis by forward and backward recursions.
In the log domain it uses the Jacobian logarithm (log-MAP); dropping the correction gives max-log-MAP,
about 0.3--0.5\,dB worse, most of which an extrinsic scaling of about 0.7 recovers.
\item Turbo codes concatenate two recursive convolutional codes through an interleaver and decode
iteratively. Recursion gives the interleaver gain; the interleaver (QPP in LTE) determines the error
floor; the constituent codes determine the waterfall. EXIT charts predict the threshold as the SNR
at which a tunnel opens between the two decoders' transfer curves.
\item LDPC codes are defined by sparse parity-check matrices and decoded by belief propagation on the
Tanner graph: sums at variable nodes, box-plus at check nodes. Min-sum with normalisation or offset
is the practical check rule; layered scheduling halves the iterations. Density evolution computes
the threshold of an ensemble, in closed form on the BEC; irregular degree distributions bring it to
within hundredths of a decibel of capacity.
\item Practical LDPC codes are protographs lifted with circulants. 5G NR uses two rate-compatible
base graphs with 51 lifting sizes; Wi-Fi, DVB-S2, 10GBASE-T, CCSDS and SSDs use QC or IRA codes.
Error floors come from trapping sets and are handled by design, decoder tweaks or an outer code.
\item Polar codes polarize $N$ channel uses into nearly perfect and nearly useless bit-channels,
freeze the bad ones and decode by successive cancellation in $O(N\log N)$. CRC-aided list decoding
makes them the best practical codes for short blocks, and they carry 5G's broadcast and control
channels.
\item The normal approximation is the yardstick at finite length: good codes sit half a decibel to a
decibel from it. HARQ with incremental redundancy uses rate-compatible codes to adapt the rate after
the fact. Spatial coupling achieves capacity with BP by threshold saturation; staircase and
soft-decision codes serve optics, and LDPC codes protect flash memory.
\end{itemize}

\begin{labbox}[Turbo, LDPC and polar codes in Python]
\raggedright
Two notebooks accompany this chapter.
\begin{itemize}
\item \lab{moderncomms-labs/labs/lab09\_ldpc\_polar.py} builds a (3,6) LDPC code by progressive edge
growth with \texttt{commlib.LDPCCode} and counts its 4-cycles; decodes it with sum-product and
normalised min-sum and plots the waterfall; compares block lengths 96, 576 and 1152; shows the
distribution of iterations with early stopping (interactive); plots channel polarization on the BEC
with \texttt{commlib.PolarCode}; and compares SC with CRC-aided SC list decoding ($L=2$ and 8, CRC-11)
at $N=256$.
\item \lab{moderncomms-labs/labs/lab23\_turbo\_exit.py} uses the new module
\texttt{commlib/turbo.py} to build the LTE turbo encoder with a QPP interleaver, verify the BCJR
decoder against brute force on a short block, plot BER iteration by iteration, compare log-MAP with
max-log-MAP, measure EXIT curves and a decoding trajectory, and run density evolution on the BEC for
regular, irregular and spatially coupled ensembles.
\end{itemize}
Every figure in this chapter is produced by \texttt{book/figscripts/ch15\_figs.py}, which also
contains batched LDPC decoders (flooding and layered; sum-product, min-sum and its corrections), a
batched CA-SCL polar decoder with Gaussian-approximation construction, and the BI-AWGN normal
approximation---good starting points for the simulation problems below.
\end{labbox}

\problems
\begin{enumerate}[label=\thechapter.\arabic*]
\item Show that for BPSK in Gaussian noise the channel LLR is $2y/\sigma^2$ and that, given $x=+1$, it
is Gaussian with variance equal to twice its mean. Show that any LLR $L$ satisfies the consistency
condition $p(L\,|\,x=+1)=e^{L}p(-L\,|\,x=+1)$, and that a Gaussian density satisfies it only if its
variance is twice its mean.
\item Prove the box-plus rule~\eqref{eq:ch15:boxplus} and the exact identity
\[L_1\boxplus L_2=\operatorname{sign}(L_1)\operatorname{sign}(L_2)\min(|L_1|,|L_2|)+\ln(1+e^{-|L_1+L_2|})-\ln(1+e^{-|L_1-L_2|}).\]
What is the largest error of the min-sum approximation, and for which inputs does it occur?
\item For the 4-state RSC with feedback $1+D+D^2$ and feed-forward $1+D^2$, draw the trellis and
find the lowest-weight input sequences that return the encoder to state 0. Explain why weight-1
inputs never do.
\item Run three steps of the BCJR algorithm by hand for the code of the previous problem with
$K=2$ information bits and two tail bits, given channel LLRs of your choice, and check the result
against brute-force enumeration of the four codewords.
\item Verify that the QPP $\pi(i)=(31i+64i^2)\bmod1024$ is a permutation, and show that it is
contention-free for 8 parallel windows of 128: for every $i<128$ the eight values
$\lfloor\pi(i+128j)/128\rfloor$, $j=0,\dots,7$, are distinct.
\item Using~\eqref{eq:ch15:floor}, estimate the error floor at $\ebno=2.5$\,dB of a rate-1/3 turbo
code with $K=6144$ whose lightest codewords are three of weight 22, each from a weight-2 input. How
much would doubling $K$ (with the same low-weight structure) lower it?
\item An LDPC ensemble has $\lambda(x)=0.25x+0.35x^2+0.4x^5$ and $\rho(x)=0.5x^6+0.5x^7$. Find its
design rate and its threshold on the BEC from~\eqref{eq:ch15:becthr}. Compare with capacity.
\item For a degree-4 check node with incoming LLRs $(+3.0,-1.2,+0.4,+2.5)$, compute all four outgoing
messages with the tanh rule, with min-sum, and with normalised min-sum ($\alpha=0.8$). Which
outgoing message is most affected by the approximation, and why?
\item Show that a 4-cycle in a QC-LDPC code lifted from shifts $s_{ij}$ exists if and only if
$s_{i_1j_1}-s_{i_1j_2}+s_{i_2j_2}-s_{i_2j_1}\equiv0\pmod Z$. Find all 4-cycles of the lifted matrix in
Figure~\ref{fig:ch15_lifting}.
\item A 5G PUSCH transport block has $A=20\,000$ bits at $R=0.8$. Determine the base graph, the CRC
lengths, the number of code blocks, $K_b$, the lifting size and the number of filler bits per code
block (TS 38.212, Section 5.2.2).
\item For the BEC with $\varepsilon=0.3$ and $N=8$, compute all eight bit-channel erasure
probabilities. Choose the best four for a $K=4$ code and compute the code's block erasure
probability under SC decoding (union bound).
\item The length-4 polar code of the worked example receives the LLRs
\[\vect L=(-0.5,\,+1.2,\,-1.8,\,-0.3).\]
Decode it by SC, and then by SCL with $L=2$ using the path
metric~\eqref{eq:ch15:pm}. Do the two decoders agree?
\item \emph{(Simulation)} Using the module \texttt{commlib.turbo}, reproduce Figure~\ref{fig:ch15_turbo_iters}
for $K=512$, and plot the average number of iterations needed when decoding stops as soon as a CRC-24
passes. Repeat with max-log-MAP and an extrinsic scaling factor swept from 0.5 to 1.
\item \emph{(Simulation)} Implement the Gaussian-approximation density evolution~\eqref{eq:ch15:ga}
for $(d_v,d_c)$-regular ensembles and find the BI-AWGN thresholds of $(3,6)$, $(4,8)$ and $(3,4)$.
Compare with Figure~\ref{fig:ch15_de_awgn} and with the exact value 0.881 for $(3,6)$.
\item \emph{(Simulation)} Using the batched polar decoder in \texttt{ch15\_figs.py}, compare the
FER of $(512,256)$ polar codes built with the Bhattacharyya construction at three design SNRs and with
the GA construction. How sensitive is CA-SCL ($L=8$) to the design SNR?
\end{enumerate}

\furtherreading
\begin{itemize}
\item T.~Richardson and R.~Urbanke, \emph{Modern Coding Theory}, Cambridge University Press, 2008.
The definitive treatment of LDPC and turbo codes, density evolution, EXIT charts and the BEC; rigorous
and beautifully organised.
\item W.~E.~Ryan and S.~Lin, \emph{Channel Codes: Classical and Modern}, Cambridge University Press,
2009. An engineering text covering BCJR, turbo, LDPC (including QC and protograph designs) and
iterative decoding in detail.
\item C.~Berrou, A.~Glavieux and P.~Thitimajshima, ``Near Shannon limit error-correcting coding and
decoding: Turbo-codes,'' \emph{Proc.\ IEEE ICC}, Geneva, 1993, pp.~1064--1070. The paper that started it.
\item L.~R.~Bahl, J.~Cocke, F.~Jelinek and J.~Raviv, ``Optimal decoding of linear codes for minimizing
symbol error rate,'' \emph{IEEE Trans.\ Information Theory}, vol.~20, no.~2, 1974.
\item R.~G.~Gallager, \emph{Low-Density Parity-Check Codes}, MIT Press, 1963; and D.~J.~C.~MacKay,
``Good error-correcting codes based on very sparse matrices,'' \emph{IEEE Trans.\ Information Theory},
vol.~45, no.~2, 1999, the full account of the rediscovery.
\item S.~ten Brink, ``Convergence behavior of iteratively decoded parallel concatenated codes,''
\emph{IEEE Trans.\ Communications}, vol.~49, no.~10, 2001. The EXIT chart.
\item T.~Richardson, A.~Shokrollahi and R.~Urbanke, ``Design of capacity-approaching irregular
low-density parity-check codes,'' \emph{IEEE Trans.\ Information Theory}, vol.~47, no.~2, 2001.
\item E.~Ar\i kan, ``Channel polarization: a method for constructing capacity-achieving codes for
symmetric binary-input memoryless channels,'' \emph{IEEE Trans.\ Information Theory}, vol.~55, no.~7,
2009; and I.~Tal and A.~Vardy, ``List decoding of polar codes,'' \emph{IEEE Trans.\ Information
Theory}, vol.~61, no.~5, 2015.
\item T.~Richardson and S.~Kudekar, ``Design of low-density parity check codes for 5G New Radio,''
\emph{IEEE Communications Magazine}, vol.~56, no.~3, 2018; and V.~Bioglio, C.~Condo and I.~Land,
``Design of polar codes in 5G New Radio,'' \emph{IEEE Communications Surveys \& Tutorials}, vol.~23,
no.~1, 2021. Readable accounts of the two NR codes by people involved in their design.
\item 3GPP TS 38.212, \emph{NR; Multiplexing and channel coding}, and TS 36.212 for LTE. The
authoritative definitions of the NR LDPC and polar codes and the LTE turbo code, with rate matching.
\item S.~Kudekar, T.~Richardson and R.~Urbanke, ``Threshold saturation via spatial coupling: why
convolutional LDPC ensembles perform so well over the BEC,'' \emph{IEEE Trans.\ Information
Theory}, vol.~57, no.~2, 2011.
\item Y.~Polyanskiy, H.~V.~Poor and S.~Verd\'u, ``Channel coding rate in the finite blocklength
regime,'' \emph{IEEE Trans.\ Information Theory}, vol.~56, no.~5, 2010. The normal approximation used
as the yardstick in Section~\ref{sec:ch15:compare}.
\end{itemize}
